%% toc: P3 35 35.5,35.6,35.7
\section{Infinitary Languages, Karp's Theorem and Scott's Theorem after Marker's Lectures}
\label{sec:p03-w03-p03-039-marker}

This section follows one presentation of the classical theory of
\(\mathcal L_{\omega_1,\omega}\): the Fall~2013 lecture notes of David Marker,
\emph{Lectures on Infinitary Model Theory}, in the January~2014 version
posted on Marker's course page \cite{Marker2014InfinitaryNotes}. These
notes are an earlier version of the book published in 2016
\cite{Marker2016Infinitary}; every locator below refers to the 2014 notes,
and nothing here is attributed to the printed book. Many results of the
notes are posed as exercises; each one used here is proved in full in the
annex, and the text says where the notes leave a proof to the reader.

The questions are these. Which formulas does an infinitary language admit,
and what does it lose with respect to first-order logic? When do two
structures satisfy the same infinitary sentences? When does a single
countable sentence pin down a countable structure up to isomorphism? And,
since the compactness theorem fails, by what argument can models still be
built? The inputs are a vocabulary \(\tau\) and \(\tau\)-structures, written
\(\mathcal M,\mathcal N\) with universes \(M,N\); the conclusions are
Karp's game and partial-isomorphism characterization of
\(\equiv_{\infty,\omega}\), Scott's isomorphism theorem, and the model
existence theorem for consistency properties.

A second construction of Scott sentences, through Morley's coded
approximants, is in \cref{sec:p03-morley-scott}; the comparison is made in
\cref{sec:p03-w03-p03-039-bridge}. The basic definition of
\(L_{\omega_1\omega}\) used elsewhere in the handbook is
\cref{def:p02-lomega-one-omega}.

\subsection{Languages, fragments and the loss of compactness}
\label{sec:p03-w03-p03-039-languages}

\begin{hbdefinition}{def:p03-w03-p03-039-language}{The languages \(\mathcal L_{\infty,\omega}\) and \(\mathcal L_{\kappa,\omega}\)}
Let \(\tau\) be a vocabulary of function, relation and constant symbols, and
fix variables \(v_\alpha\), one for each ordinal \(\alpha\). The formulas of
\(\mathcal L_{\infty,\omega}(\tau)\) are generated as follows: terms and
atomic formulas are as in first-order logic; if \(\varphi\) is a formula then
so is \(\neg\varphi\); if \(X\) is a \emph{set} of formulas then
\(\bigwedge_{\varphi\in X}\varphi\) and \(\bigvee_{\varphi\in X}\varphi\) are
formulas; and if \(\varphi\) is a formula then so are \(\exists v_\alpha\,\varphi\)
and \(\forall v_\alpha\,\varphi\). Satisfaction extends the first-order clauses
by: \(\mathcal M\models\bigvee_{\varphi\in X}\varphi\) exactly when
\(\mathcal M\models\varphi\) for some \(\varphi\in X\), and
\(\mathcal M\models\bigwedge_{\varphi\in X}\varphi\) exactly when
\(\mathcal M\models\varphi\) for all \(\varphi\in X\). Binary \(\land,\lor\)
and the connectives \(\to,\leftrightarrow\) are abbreviations; here
\(\varphi\to\psi\) abbreviates \(\bigvee\{\neg\varphi,\psi\}\).

For an infinite cardinal \(\kappa\), \(\mathcal L_{\kappa,\omega}(\tau)\) uses
only the variables \(v_\alpha\) with \(\alpha<\kappa\) and only conjunctions
and disjunctions over sets \(X\) with \(\lvert X\rvert<\kappa\). Thus
\(\mathcal L_{\omega,\omega}\) is first-order logic and
\(\mathcal L_{\omega_1,\omega}\) admits countable conjunctions and
disjunctions. Every quantifier binds one variable, so quantifier strings stay
finite. Structures \(\mathcal M,\mathcal N\) satisfy
\(\mathcal M\equiv_{\infty,\omega}\mathcal N\) when they satisfy the same
\(\mathcal L_{\infty,\omega}\)-sentences; \(\equiv_{\kappa,\omega}\) is
defined in the same way with \(\mathcal L_{\kappa,\omega}\)-sentences
\cite[pp.~4--5, Def.~1.3]{Marker2014InfinitaryNotes}.
\end{hbdefinition}

A formula of \(\mathcal L_{\infty,\omega}\) may have infinitely many free
variables, for instance \(\bigwedge_{n<\omega}v_n=v_n\). A subformula of a
\emph{sentence} has only finitely many, because the free variables of a
subformula are bound by the finitely many quantifiers above it in the
formation tree
\cite[Ex.~1.1, p.~5]{Marker2014InfinitaryNotes}. This is the point at which
the present definition differs from \cref{def:p02-lomega-one-omega}, which
builds the finite-free-variable condition into the grammar; the two languages
have the same sentences.

\begin{hbtheorem}{thm:p03-w03-p03-039-iso-invariance}{proved=annex:p03-w03-p03-039-iso-invariance}{Isomorphism preserves infinitary truth}
If \(h:\mathcal M\to\mathcal N\) is an isomorphism, then for every formula
\(\varphi\) of \(\mathcal L_{\infty,\omega}\) and every assignment \(s\) into
\(M\), \(\mathcal M,s\models\varphi\) if and only if
\(\mathcal N,h\circ s\models\varphi\). In particular
\(\mathcal M\cong\mathcal N\) implies
\(\mathcal M\equiv_{\infty,\omega}\mathcal N\)
\cite[Ex.~1.4, p.~5]{Marker2014InfinitaryNotes}.
\end{hbtheorem}

\begin{hbtheorem}{thm:p03-w03-p03-039-compactness-failure}{proved=annex:p03-w03-p03-039-compactness-failure}{Compactness fails for \(\mathcal L_{\omega_1,\omega}\)}
Let \(\tau\) have constant symbols \(d,c_0,c_1,\ldots\) and nothing else, and let
\[
 \Gamma=\{d\ne c_i:i\in\omega\}\cup
        \Bigl\{\forall v\bigvee_{i\in\omega}v=c_i\Bigr\}.
\]
Every finite subset of \(\Gamma\) has a model, but \(\Gamma\) has none. Hence
the compactness theorem fails for \(\mathcal L_{\omega_1,\omega}\) and for
\(\mathcal L_{\infty,\omega}\) \cite[Ex.~1.5, p.~5]{Marker2014InfinitaryNotes}.
\end{hbtheorem}

The notes stress that this failure shapes the whole subject: models are
built below by Henkin constructions from consistency properties rather than
by compactness \cite[pp.~5--6, 39]{Marker2014InfinitaryNotes}. The upward
L\"owenheim--Skolem theorem fails as well
(\cref{thm:p03-w03-p03-039-upward-ls-failure}).

\begin{hbdefinition}{def:p03-w03-p03-039-formal-negation}{Formal negation}
For each formula \(\varphi\) of \(\mathcal L_{\infty,\omega}\) its formal
negation \({\sim}\varphi\) is defined by recursion on \(\varphi\): for atomic
\(\varphi\), \({\sim}\varphi\) is \(\neg\varphi\); \({\sim}(\neg\varphi)\) is
\(\varphi\); \({\sim}\bigwedge_{\varphi\in X}\varphi\) is
\(\bigvee_{\varphi\in X}{\sim}\varphi\) and
\({\sim}\bigvee_{\varphi\in X}\varphi\) is
\(\bigwedge_{\varphi\in X}{\sim}\varphi\); \({\sim}\exists v\,\varphi\) is
\(\forall v\,{\sim}\varphi\) and \({\sim}\forall v\,\varphi\) is
\(\exists v\,{\sim}\varphi\)
\cite[Def.~1.16, p.~8]{Marker2014InfinitaryNotes}.
\end{hbdefinition}

\begin{hblemma}{thm:p03-w03-p03-039-formal-negation}{proved=annex:p03-w03-p03-039-formal-negation}{Formal negation is negation}
For every formula \(\varphi\) of \(\mathcal L_{\infty,\omega}\), every
structure \(\mathcal M\) and every assignment \(s\),
\(\mathcal M,s\models\neg\varphi\) if and only if
\(\mathcal M,s\models{\sim}\varphi\). The formula \({\sim}\varphi\) has the
same free variables as \(\varphi\), and substituting a term for a free
variable commutes with \({\sim}\)
\cite[Ex.~1.17, p.~8]{Marker2014InfinitaryNotes}.
\end{hblemma}

\begin{hbdefinition}{def:p03-w03-p03-039-fragment}{Fragment}
A set \(A\) of \(\mathcal L_{\infty,\omega}\)-formulas is a \emph{fragment}
if there is an infinite set \(V\) of variables such that every variable
occurring in a member of \(A\) lies in \(V\), and
\begin{enumerate}
\item every atomic formula built from the symbols of \(\tau\) and variables
      from \(V\) is in \(A\);
\item \(A\) is closed under subformulas;
\item if \(\varphi\in A\), \(v\) is free in \(\varphi\) and \(t\) is a term
      whose variables lie in \(V\), the result of substituting \(t\) for the
      free occurrences of \(v\) is in \(A\);
\item \(A\) is closed under \({\sim}\);
\item \(A\) is closed under \(\neg\), binary \(\land\) and \(\lor\), and
      \(\exists v\), \(\forall v\) for \(v\in V\).
\end{enumerate}
For a fragment \(A\), \(\mathcal M\equiv_A\mathcal N\) means agreement on
the sentences of \(A\), and \(\mathcal N\prec_A\mathcal M\) means that
\(\mathcal N\) is a substructure of \(\mathcal M\) and, for every
\(\varphi(\bar v)\in A\) and every tuple \(\bar a\) from \(N\),
\(\mathcal N\models\varphi(\bar a)\) if and only if
\(\mathcal M\models\varphi(\bar a)\)
\cite[Def.~1.18, p.~9]{Marker2014InfinitaryNotes}.
\end{hbdefinition}

The notes list the atomic formulas of clause~1 as those using the constant
symbols of \(\tau\) and variables from \(V\); the atomic formulas meant are
those of the whole vocabulary, and the counting below depends on that.

\begin{hblemma}{thm:p03-w03-p03-039-countable-fragment}{proved=annex:p03-w03-p03-039-countable-fragment}{Countable sets of sentences lie in countable fragments}
Let \(\tau\) be countable and let \(T\) be a countable set of
\(\mathcal L_{\omega_1,\omega}(\tau)\)-formulas. Let \(V\) be a countably
infinite set of variables containing every variable that occurs in a member
of \(T\). There is a countable fragment \(A\subseteq\mathcal L_{\omega_1,\omega}(\tau)\)
with witness \(V\) and \(T\subseteq A\), and there is a least such fragment:
it is contained in every fragment with witness \(V\) that contains \(T\)
\cite[Ex.~1.19, p.~9]{Marker2014InfinitaryNotes}.
\end{hblemma}

\begin{hbwarning}{def:p03-w03-p03-039-fragment-size}{The size of a fragment is bounded below by the vocabulary}
The notes pose the general form of \cref{thm:p03-w03-p03-039-countable-fragment}
for a regular \(\kappa\) and a set \(T\) of fewer than \(\kappa\)
\(\mathcal L_{\kappa,\omega}\)-sentences, concluding \(\lvert A\rvert<\kappa\)
\cite[Ex.~1.19(a), p.~9]{Marker2014InfinitaryNotes}. That conclusion needs
\(\lvert\tau\rvert<\kappa\) as well: clause~1 of
\cref{def:p03-w03-p03-039-fragment} puts an atomic formula
\(R(v,\ldots,v)\) in \(A\) for every relation symbol \(R\), and an atomic
formula \(f(v,\ldots,v)=v\) for every function symbol \(f\), so
\(\lvert A\rvert\ge\lvert\tau\rvert\). With \(T=\{\forall v\,v=v\}\) and
\(\lvert\tau\rvert=\kappa\) no fragment containing \(T\) has fewer than
\(\kappa\) members. The minimality claim is relative to the witness \(V\):
fragments with different variable sets contain different atomic formulas and
need not be comparable.
\end{hbwarning}

\begin{hbtheorem}{thm:p03-w03-p03-039-downward-ls}{proved=annex:p03-w03-p03-039-downward-ls}{Downward L\"owenheim--Skolem for fragments}
Let \(A\) be a fragment of \(\mathcal L_{\infty,\omega}(\tau)\) in which every
formula has only finitely many free variables. Let \(\mathcal M\) be a
\(\tau\)-structure and \(X\subseteq M\). There is
\(\mathcal N\prec_A\mathcal M\) with \(X\subseteq N\) and
\(\lvert N\rvert\le\max(\lvert A\rvert,\lvert X\rvert)\). In particular, if
\(\tau\) is countable and \(A\) is a countable fragment of
\(\mathcal L_{\omega_1,\omega}(\tau)\), every \(\tau\)-structure has a
countable \(A\)-elementary substructure
\cite[Thm.~1.23, p.~10]{Marker2014InfinitaryNotes}.
\end{hbtheorem}

The notes outline the proof by expanding \(\mathcal M\) with Skolem functions
for the formulas of \(A\) and observing that every substructure of the
expansion is \(A\)-elementary \cite[Ex.~1.24, p.~10]{Marker2014InfinitaryNotes}.
The annex takes the same Skolem hull without naming the new function
symbols. The finite-free-variable hypothesis is what makes each Skolem
function finitary. Without a countable fragment the conclusion can fail: the
notes ask for a countable vocabulary and an \(\mathcal L_{\omega_1,\omega}\)
theory all of whose models have cardinality at least \(2^{\aleph_0}\)
\cite[Ex.~1.15, p.~8]{Marker2014InfinitaryNotes}; such a theory has
uncountably many sentences, and the theorem then yields only submodels of
size \(\lvert A\rvert\).

\begin{hbtheorem}{thm:p03-w03-p03-039-chang}{proved=annex:p03-w03-p03-039-chang}{Chang's reduction to first-order omitting types}
Let \(\tau\) be countable and let \(T\) be a countable set of
\(\mathcal L_{\omega_1,\omega}(\tau)\)-sentences. There are a countable
vocabulary \(\tau^*\supseteq\tau\), a first-order \(\tau^*\)-theory \(T^*\) and
a countable set \(\Gamma\) of partial first-order \(\tau^*\)-types such that
\begin{enumerate}
\item if \(\mathcal M\models T^*\) and \(\mathcal M\) omits every type in
      \(\Gamma\), then the \(\tau\)-reduct of \(\mathcal M\) is a model of \(T\);
\item every model of \(T\) has a \(\tau^*\)-expansion that is a model of \(T^*\)
      omitting every type in \(\Gamma\)
\end{enumerate}
\cite[Thm.~1.31 and Ex.~1.32, pp.~11--12]{Marker2014InfinitaryNotes}.
\end{hbtheorem}

\begin{hbmethod}{def:p03-w03-p03-039-chang-construction}{The Chang expansion}
Let \(A\) be the least countable fragment containing \(T\)
(\cref{thm:p03-w03-p03-039-countable-fragment}). For each \(\varphi\in A\),
fix the increasing list \(\bar x_\varphi\) of its free variables and add a
relation symbol \(R_\varphi\) of that arity. \(T^*\) consists of:
\begin{enumerate}
\item \(\forall\bar x_\varphi(R_\varphi\leftrightarrow\varphi)\) for atomic \(\varphi\);
\item \(\forall\bar x_\varphi(R_\varphi\leftrightarrow\neg R_\theta)\) when
      \(\varphi\) is \(\neg\theta\);
\item \(\forall\bar x_\varphi(R_\varphi\to R_\theta)\) for each \(\theta\in X\),
      when \(\varphi\) is \(\bigwedge_{\theta\in X}\theta\), with the type
      \(\gamma_\varphi=\{\neg R_\varphi\}\cup\{R_\theta:\theta\in X\}\);
\item \(\forall\bar x_\varphi(R_\theta\to R_\varphi)\) for each \(\theta\in X\),
      when \(\varphi\) is \(\bigvee_{\theta\in X}\theta\), with the type
      \(\gamma_\varphi=\{R_\varphi\}\cup\{\neg R_\theta:\theta\in X\}\);
\item \(\forall\bar x_\varphi(R_\varphi\leftrightarrow\exists w\,R_\theta)\) when
      \(\varphi\) is \(\exists w\,\theta\), and
      \(\forall\bar x_\varphi(R_\varphi\leftrightarrow\forall w\,R_\theta)\) when
      \(\varphi\) is \(\forall w\,\theta\);
\item \(R_\varphi\) for each sentence \(\varphi\in T\).
\end{enumerate}
In each clause \(R_\theta\) is applied to the free variables of \(\theta\),
which lie among \(\bar x_\varphi\) and, in the quantifier clauses, \(w\).
\(\Gamma\) is the set of all \(\gamma_\varphi\). The printed universal clause
reads \(\forall\bar v(R_\varphi\leftrightarrow\forall\bar v\forall w\,R_\theta)\)
\cite[p.~12]{Marker2014InfinitaryNotes}; the inner \(\forall\bar v\) is a
misprint and is dropped here, since with it the clause would not track the
satisfaction of \(\forall w\,\theta\) at a given tuple.
\end{hbmethod}

\subsection{Size of formulas, ultraproducts, axiomatizable and pseudo-elementary classes}
\label{sec:p03-w03-p03-039-expressive}

This subsection collects the remaining results of the notes' first chapter
that concern the logic itself rather than fragments: how many subformulas a
sentence has, how ultraproducts interact with countable conjunctions, which
familiar classes are axiomatizable, the failure of the upward
L\"owenheim--Skolem theorem, the separation of \(\equiv_{\kappa,\omega}\) from
\(\equiv_{\infty,\omega}\), and pseudo-elementary classes. The notes pose all of
them as exercises; every one is proved in the annex.

\begin{hbtheorem}{thm:p03-w03-p03-039-subformula-count}{proved=annex:p03-w03-p03-039-subformula-count}{Subformula counts and regularity}
Let \(\kappa\) be an infinite cardinal.
\begin{enumerate}
\item If \(\kappa\) is regular, every formula of \(\mathcal L_{\kappa,\omega}(\tau)\)
      has fewer than \(\kappa\) subformulas.
\item If \(\kappa\) is singular, there is a sentence of
      \(\mathcal L_{\kappa,\omega}\) in the empty vocabulary with \(\kappa\)
      distinct subformulas.
\end{enumerate}
\cite[Ex.~1.2, p.~5]{Marker2014InfinitaryNotes}.
\end{hbtheorem}

The notes give this as one reason for restricting attention to
\(\mathcal L_{\kappa,\omega}\) with \(\kappa\) regular
\cite[Ex.~1.2, p.~5]{Marker2014InfinitaryNotes}; the singular case is the one
that Nadel's chapter also sets aside (\cref{def:p03-w03-p03-039-nadel-logics}).

\begin{hbdefinition}{def:p03-w03-p03-039-ultraproduct}{Ultraproducts and \(\sigma\)-complete ultrafilters}
Let \((\mathcal M_i:i\in I)\) be \(\tau\)-structures and \(U\) an ultrafilter
on \(I\). For \(f,g\in\prod_{i\in I}M_i\) put \(f\sim_U g\) when
\(\{i:f(i)=g(i)\}\in U\), and write \([f]\) for the class of \(f\). The
ultraproduct \(\prod_{i\in I}\mathcal M_i/U\) has universe the set of classes;
a constant \(c\) is interpreted as \([i\mapsto c^{\mathcal M_i}]\), a function
symbol \(F\) by
\(F([f_1],\dots,[f_n])=[i\mapsto F^{\mathcal M_i}(f_1(i),\dots,f_n(i))]\), and
a relation symbol \(R\) holds of \(([f_1],\dots,[f_n])\) when
\(\{i:\mathcal M_i\models R(f_1(i),\dots,f_n(i))\}\in U\). The ultrafilter
\(U\) is \emph{\(\sigma\)-complete} if the intersection of any countably many
members of \(U\) is in \(U\).
\end{hbdefinition}

\begin{hbtheorem}{thm:p03-w03-p03-039-ultraproducts}{proved=annex:p03-w03-p03-039-ultraproducts}{Ultraproducts and \(\mathcal L_{\omega_1,\omega}\)}
\begin{enumerate}
\item Let \(\tau=\{c_n:n\in\omega\}\), let every \(\mathcal M_i\)
      \((i\in\omega)\) be \(\omega\) with \(c_n\) interpreted as \(n\), and let
      \(\varphi\) be \(\forall v\bigvee_{n\in\omega}v=c_n\). Then
      \(\mathcal M_i\models\varphi\) for every \(i\), and
      \(\prod_{i\in\omega}\mathcal M_i/U\models\neg\varphi\) for every
      nonprincipal ultrafilter \(U\) on \(\omega\).
\item Let \(U\) be a \(\sigma\)-complete ultrafilter on \(I\). For every
      formula \(\varphi(v_1,\dots,v_n)\) of \(\mathcal L_{\omega_1,\omega}(\tau)\)
      with finitely many free variables and all \(f_1,\dots,f_n\in\prod_iM_i\),
      \[
      \textstyle\prod_i\mathcal M_i/U\models\varphi([f_1],\dots,[f_n])
      \iff\{i\in I:\mathcal M_i\models\varphi(f_1(i),\dots,f_n(i))\}\in U .
      \]
\end{enumerate}
\cite[Ex.~1.6, p.~6]{Marker2014InfinitaryNotes}.
\end{hbtheorem}

Part~1 is the ultraproduct form of the failure of compactness: \L o\'s's
theorem breaks down at the countable disjunction, because the sets
\(\{i:f(i)=n\}\) may all lie outside \(U\) while their union is \(I\). Part~2
shows that countable completeness of the ultrafilter is exactly what repairs
it.

\begin{hbtheorem}{thm:p03-w03-p03-039-axiomatizable}{proved=annex:p03-w03-p03-039-axiomatizable}{Classes axiomatizable in \(\mathcal L_{\omega_1,\omega}\)}
Each of the following classes is the class of models of a set of
\(\mathcal L_{\omega_1,\omega}\)-sentences in the vocabulary indicated; in
cases i)--vi) and ix) (when \(T\) is countable) a single sentence suffices,
and in vii) and viii) the countable conjunction of the set does:
\begin{enumerate}
\item[i)] torsion abelian groups, in \(\{+,-,0\}\);
\item[ii)] finitely generated groups, in \(\{\cdot,{}^{-1},1\}\);
\item[iii)] groups that are not finitely generated, in \(\{\cdot,{}^{-1},1\}\);
\item[iv)] linear orders isomorphic to \((\mathbb Z,<)\), in \(\{<\}\);
\item[v)] archimedean ordered fields, in \(\{+,\cdot,-,0,1,<\}\);
\item[vi)] connected graphs, in \(\{E\}\);
\item[vii)] recursively saturated models of Peano arithmetic, in
      \(\{+,\cdot,0,1,<\}\);
\item[viii)] \(\omega\)-models of \(\mathrm{ZFC}\), in \(\{\in\}\);
\item[ix)] models of a first-order theory \(T\) omitting a type \(p(\bar v)\),
      in a countable vocabulary
\end{enumerate}
\cite[Ex.~1.7, p.~6]{Marker2014InfinitaryNotes}.
\end{hbtheorem}

The notes leave the meaning of two items implicit. A model \(\mathcal M\) of
\(\mathrm{PA}\) is \emph{recursively saturated} when for every recursive set
\(p(v,w_1,\dots,w_k)\) of first-order formulas, with free variables among
\(v,w_1,\dots,w_k\), and every \(\bar a\in M^k\), if every finite subset of
\(p(v,\bar a)\) is satisfied in \(\mathcal M\) then some element satisfies all
of \(p(v,\bar a)\). A model of \(\mathrm{ZFC}\) is an \emph{\(\omega\)-model}
when every element that it takes to be a natural number is the value of a
numeral: the values \(0^{\mathcal M}=\varnothing^{\mathcal M}\) and
\((n+1)^{\mathcal M}=n^{\mathcal M}\cup\{n^{\mathcal M}\}\) exhaust
\(\omega^{\mathcal M}\). Items ii) and ix) overlap
\cref{thm:p03-w03-p03-039-nadel-expressive}, where Nadel's versions of the
same sentences are proved; the annex reuses those sentences and proves the
remaining items. The notes add that none of the nine classes is first-order
axiomatizable \cite[Ex.~1.8, p.~6]{Marker2014InfinitaryNotes}.

\begin{hbtheorem}{thm:p03-w03-p03-039-upward-ls-failure}{proved=annex:p03-w03-p03-039-upward-ls-failure}{The upward L\"owenheim--Skolem theorem fails}
\begin{enumerate}
\item The sentence of \cref{thm:p03-w03-p03-039-axiomatizable}~iv) has an
      infinite model, and every model of it is countable.
\item Every finitely generated group is countable, and the infinite cyclic
      group is an infinite finitely generated group.
\item Every archimedean ordered field has cardinality at most
      \(2^{\aleph_0}\); \(\mathbb Q\) and \(\mathbb R\) are archimedean ordered
      fields.
\end{enumerate}
Hence there are sentences of \(\mathcal L_{\omega_1,\omega}\) with infinite
models and no model of some larger cardinality
\cite[p.~7]{Marker2014InfinitaryNotes}.
\end{hbtheorem}

The notes justify item~3 by saying that an archimedean ordered field is
isomorphic to a subfield of \(\mathbb R\) \cite[p.~7]{Marker2014InfinitaryNotes};
the annex proves the part of that statement the bound needs, an
order-preserving injection into \(\mathbb R\).

The notes next ask for sentences \(\Phi_\alpha\) of \(\mathcal L_{\infty,\omega}\)
describing \((\alpha,<)\) up to isomorphism, by induction on \(\alpha\)
\cite[Ex.~1.9, p.~7]{Marker2014InfinitaryNotes}. This is the exercise
Nadel poses on p.~274 of his chapter, and it is proved as
\cref{thm:p03-w03-p03-039-nadel-ordinals}: the sentence \(\sigma_\alpha\)
there is a \(\Phi_\alpha\) for every \(\alpha\ge1\).

\begin{hblemma}{thm:p03-w03-p03-039-formula-count}{proved=annex:p03-w03-p03-039-formula-count}{Counting formulas of \(\mathcal L_{\kappa,\omega}\)}
Let \(\kappa\) be an infinite regular cardinal and \(\lvert\tau\rvert\le\kappa\).
The formulas of \(\mathcal L_{\kappa,\omega}(\tau)\) form a set of cardinality
at most \(2^\kappa\).
\end{hblemma}

\begin{hbtheorem}{thm:p03-w03-p03-039-ordinal-equivalence}{proved=annex:p03-w03-p03-039-ordinal-equivalence}{Ordinals that \(\mathcal L_{\kappa,\omega}\) cannot tell apart}
Let \(\kappa\) be an infinite regular cardinal. There are ordinals
\(\alpha\ne\beta\), both nonzero and below \((2^{2^\kappa})^+\), with
\((\alpha,<)\equiv_{\kappa,\omega}(\beta,<)\). Since
\((\alpha,<)\not\equiv_{\infty,\omega}(\beta,<)\) by
\cref{thm:p03-w03-p03-039-nadel-ordinals}, for every regular \(\kappa\) there
are structures that are \(\equiv_{\kappa,\omega}\)-equivalent and not
\(\equiv_{\infty,\omega}\)-equivalent
\cite[Ex.~1.10 and p.~7]{Marker2014InfinitaryNotes}.
\end{hbtheorem}

\begin{hbdefinition}{def:p03-w03-p03-039-pc}{\(\mathrm{PC}_{\omega_1,\omega}\)-, \(\mathrm{PC}\)- and \(\mathrm{PC}_\delta\)-classes}
A class \(\mathcal K\) of \(\tau\)-structures is a
\emph{\(\mathrm{PC}_{\omega_1,\omega}\)-class} if there are a vocabulary
\(\tau^*\supseteq\tau\) and \(\varphi\in\mathcal L_{\omega_1,\omega}(\tau^*)\)
such that \(\mathcal K\) is the class of \(\tau\)-reducts of models of
\(\varphi\). It is a \emph{\(\mathrm{PC}\)-class} if \(\varphi\) can be taken
first-order, and a \emph{\(\mathrm{PC}_\delta\)-class} if \(\varphi\) can be
replaced by a first-order theory
\cite[Def.~1.11, p.~7]{Marker2014InfinitaryNotes}.
\end{hbdefinition}

\begin{hbtheorem}{thm:p03-w03-p03-039-pc-examples}{proved=annex:p03-w03-p03-039-pc-examples}{Six \(\mathrm{PC}_{\omega_1,\omega}\)-classes}
The following classes are \(\mathrm{PC}_{\omega_1,\omega}\)-classes; a), c),
e) and f) are even \(\mathrm{PC}\)-classes:
\begin{enumerate}
\item[a)] orderable groups, that is, groups admitting a linear order \(<\)
      with \(x<y\Rightarrow zx<zy\wedge xz<yz\);
\item[b)] free groups;
\item[c)] \(1\)-transitive linear orders, that is, linear orders in which
      for all \(a,b\) some order automorphism takes \(a\) to \(b\);
\item[d)] fields that are purely transcendental extensions of \(\mathbb Q\);
\item[e)] ordered fields that are not Dedekind complete, that is, that have a
      nonempty subset bounded above with no least upper bound;
\item[f)] ordered fields with an integer part: an ordered subring \(Z\) with
      no element strictly between \(0\) and \(1\) such that every element of
      the field lies within distance at most \(1\) of an element of \(Z\)
\end{enumerate}
\cite[Ex.~1.12, p.~7]{Marker2014InfinitaryNotes}.
\end{hbtheorem}

In b) a group is free on \(B\subseteq G\) when every element of \(G\) is the
value of exactly one reduced word over \(B\), a reduced word being a finite
sequence of letters \(b^{\pm1}\) (\(b\in B\)) with no letter adjacent to its
inverse; in d) ``purely transcendental'' means \(K=\mathbb Q(B)\) for a set
\(B\) algebraically independent over \(\mathbb Q\), possibly empty.

\begin{hbtheorem}{thm:p03-w03-p03-039-silver}{proved=annex:p03-w03-p03-039-silver}{Silver's class}
Let \(\tau=\{U\}\) with \(U\) unary and
\(\mathcal K=\{\mathcal M:\lvert M\rvert\le2^{\lvert U(\mathcal M)\rvert}
\text{ and }\lvert M\setminus U(\mathcal M)\rvert=\lvert M\rvert\}\).
\begin{enumerate}
\item \(\mathcal K\) is a \(\mathrm{PC}\)-class.
\item For an infinite cardinal \(\kappa\), \(\mathcal K\) has exactly one model of
      cardinality \(\kappa\) up to isomorphism if and only if \(\kappa\) is a
      strong limit cardinal, that is, if and only if \(\kappa=\beth_\alpha\)
      with \(\alpha=0\) or \(\alpha\) a limit ordinal.
\end{enumerate}
\cite[Ex.~1.13, pp.~7--8]{Marker2014InfinitaryNotes}.
\end{hbtheorem}

The notes state part~2 as ``\(\kappa=\beth_\alpha\) for some limit ordinal
\(\alpha\)'' \cite[p.~8]{Marker2014InfinitaryNotes}; this is the statement
above when \(0\) is counted among the limit ordinals, and the case
\(\kappa=\aleph_0=\beth_0\) is a case of categoricity. The example is the
notes' witness that Morley's categoricity theorem does not extend naively to
\(\mathrm{PC}\)-classes: \(\mathcal K\) is categorical in \(\beth_\omega\) and
not in \(\aleph_1\), because \(\aleph_1\le2^{\aleph_0}\) admits the two
values \(\lvert U\rvert=\aleph_0\) and \(\lvert U\rvert=\aleph_1\).

\begin{hbtheorem}{thm:p03-w03-p03-039-large-models}{proved=annex:p03-w03-p03-039-large-models}{A countable vocabulary with only large models}
Let \(\tau=\{U_n:n\in\omega\}\) with each \(U_n\) unary, and for
\(S\subseteq\omega\) let
\(\theta_S=\exists x\bigl(\bigwedge_{n\in S}U_n(x)\wedge\bigwedge_{n\notin S}\neg U_n(x)\bigr)\).
The \(\mathcal L_{\omega_1,\omega}\)-theory \(T=\{\theta_S:S\subseteq\omega\}\)
has a model, and every model of \(T\) has cardinality at least \(2^{\aleph_0}\)
\cite[Ex.~1.15, p.~8]{Marker2014InfinitaryNotes}.
\end{hbtheorem}

This is the example promised after \cref{thm:p03-w03-p03-039-downward-ls}:
\(T\) has \(2^{\aleph_0}\) sentences, so no countable fragment contains it.

\subsection{Partial isomorphisms and Karp's theorem}
\label{sec:p03-w03-p03-039-karp}

\begin{hbdefinition}{def:p03-w03-p03-039-partial-iso}{Partial isomorphism systems and the game \(G(\mathcal M,\mathcal N)\)}
A \emph{partial \(\tau\)-embedding} from \(\mathcal M\) to \(\mathcal N\) is a
map \(f:D\to N\) with \(D\subseteq M\) such that for every atomic formula
\(\varphi(\bar v)\) and every tuple \(\bar a\) from \(D\),
\(\mathcal M\models\varphi(\bar a)\) if and only if
\(\mathcal N\models\varphi(f(\bar a))\); such a map is injective, because
equality is atomic. A \emph{partial isomorphism system} between
\(\mathcal M\) and \(\mathcal N\) is a nonempty set \(P\) of partial
\(\tau\)-embeddings such that
\begin{enumerate}
\item for every \(f\in P\) and \(a\in M\) there is \(g\in P\) with
      \(g\supseteq f\) and \(a\in\operatorname{dom}(g)\);
\item for every \(f\in P\) and \(b\in N\) there is \(g\in P\) with
      \(g\supseteq f\) and \(b\in\operatorname{img}(g)\).
\end{enumerate}
Write \(\mathcal M\cong_p\mathcal N\) if such a system exists. In the game
\(G(\mathcal M,\mathcal N)\), at stage \(n\) Player~I plays either
\(a_n\in M\) or \(b_n\in N\), and Player~II answers with \(b_n\in N\) or
\(a_n\in M\) respectively; Player~II wins the play if the map
\(a_n\mapsto b_n\) is a partial \(\tau\)-embedding
\cite[Defs.~2.1 and~2.3, p.~13]{Marker2014InfinitaryNotes}.
\end{hbdefinition}

In clause~2 the notes print ``for all \(g\in P\) \ldots there is \(g\in P\)
such that \(g\supseteq f\)''; the clause is read with the evident symmetry to
clause~1, as above.

\begin{hblemma}{thm:p03-w03-p03-039-countable-back-forth}{proved=annex:p03-w03-p03-039-countable-back-forth}{Back and forth between countable structures}
If \(\mathcal M\) and \(\mathcal N\) are at most countable and
\(\mathcal M\cong_p\mathcal N\), then \(\mathcal M\cong\mathcal N\)
\cite[Ex.~2.2, p.~13]{Marker2014InfinitaryNotes}.
\end{hblemma}

\begin{hbtheorem}{thm:p03-w03-p03-039-karp}{proved=annex:p03-w03-p03-039-karp}{Karp's theorem}
For \(\tau\)-structures \(\mathcal M\) and \(\mathcal N\) the following are
equivalent:
\begin{enumerate}
\item \(\mathcal M\equiv_{\infty,\omega}\mathcal N\);
\item \(\mathcal M\cong_p\mathcal N\);
\item there is a partial isomorphism system between \(\mathcal M\) and
      \(\mathcal N\) all of whose members have finite domain;
\item Player~II has a winning strategy in \(G(\mathcal M,\mathcal N)\)
\end{enumerate}
\cite[Thm.~2.4, pp.~13--14]{Marker2014InfinitaryNotes}.
\end{hbtheorem}

No cardinality hypothesis enters \cref{thm:p03-w03-p03-039-karp}; the
countability hypothesis enters only when \(\cong_p\) is turned into
\(\cong\) by \cref{thm:p03-w03-p03-039-countable-back-forth}. The next
example shows that the hypothesis cannot be dropped.

\begin{hbexample}{def:p03-w03-p03-039-q-r}{\((\mathbb Q,<)\) and \((\mathbb R,<)\)}
Let \(P\) be the set of finite order-preserving maps from a finite subset of
\(\mathbb Q\) into \(\mathbb R\). The empty map is in \(P\). Given
\(f\in P\) with domain \(q_1<\cdots<q_k\) and a new \(q\in\mathbb Q\), the
point \(q\) lies in one of the intervals \((-\infty,q_1)\),
\((q_i,q_{i+1})\), \((q_k,\infty)\) (or \(\mathbb Q\) when \(k=0\)), and the
corresponding interval of \(\mathbb R\) determined by the images is nonempty
because \(\mathbb R\) is dense and has no endpoints; sending \(q\) to any
point of it extends \(f\) inside \(P\). The same argument with the roles
exchanged, using that \(\mathbb Q\) is dense without endpoints, extends the
image. Order-preserving maps between linear orders preserve and reflect
\(<\) and \(=\), so they are partial embeddings, and \(P\) is a partial
isomorphism system. By \cref{thm:p03-w03-p03-039-karp},
\((\mathbb Q,<)\equiv_{\infty,\omega}(\mathbb R,<)\), although the two
orders are not isomorphic, having different cardinalities. Consequently
the Scott sentence of \((\mathbb Q,<)\) defined below holds in
\((\mathbb R,<)\): Scott's theorem characterizes a countable structure only
among countable structures.
\end{hbexample}

\begin{hbdefinition}{def:p03-w03-p03-039-approximants}{The relations \(\sim_\alpha\), quantifier rank and \(\equiv_\alpha\)}
For \(\tau\)-structures \(\mathcal M,\mathcal N\) (nonempty, as always),
tuples \(\bar a\in M^n\), \(\bar b\in N^n\), \(n=0,1,2,\ldots\), and ordinals
\(\alpha\), define
\begin{enumerate}
\item \((\mathcal M,\bar a)\sim_0(\mathcal N,\bar b)\) if \(\bar a\) and
      \(\bar b\) satisfy the same atomic \(\tau\)-formulas;
\item \((\mathcal M,\bar a)\sim_{\alpha+1}(\mathcal N,\bar b)\) if for every
      \(c\in M\) there is \(d\in N\) with
      \((\mathcal M,\bar a,c)\sim_\alpha(\mathcal N,\bar b,d)\), and for every
      \(d\in N\) there is \(c\in M\) with
      \((\mathcal M,\bar a,c)\sim_\alpha(\mathcal N,\bar b,d)\);
\item for a limit \(\beta\), \((\mathcal M,\bar a)\sim_\beta(\mathcal N,\bar b)\)
      if \((\mathcal M,\bar a)\sim_\alpha(\mathcal N,\bar b)\) for all
      \(\alpha<\beta\).
\end{enumerate}
The quantifier rank \(\operatorname{qr}\) is \(0\) on atomic formulas,
\(\operatorname{qr}(\neg\varphi)=\operatorname{qr}(\varphi)\),
\(\operatorname{qr}(\bigwedge_{\varphi\in X}\varphi)=
\operatorname{qr}(\bigvee_{\varphi\in X}\varphi)=
\sup_{\varphi\in X}\operatorname{qr}(\varphi)\), and
\(\operatorname{qr}(\exists v\,\varphi)=\operatorname{qr}(\forall v\,\varphi)=
\operatorname{qr}(\varphi)+1\). \(\mathcal M\equiv_\alpha\mathcal N\) means
that \(\mathcal M\) and \(\mathcal N\) satisfy the same
\(\mathcal L_{\infty,\omega}\)-sentences of quantifier rank at most \(\alpha\);
\(\mathcal M\sim_\alpha\mathcal N\) is the case \(n=0\)
\cite[pp.~15--16, Defs.~2.9 and~2.10]{Marker2014InfinitaryNotes}.
\end{hbdefinition}

The successor clause does not mention \(\sim_\alpha\) itself; the next lemma
shows that it follows.

\begin{hblemma}{thm:p03-w03-p03-039-approximant-basics}{proved=annex:p03-w03-p03-039-approximant-basics}{Monotonicity and equivalence}
For all ordinals \(\alpha\le\beta\),
\((\mathcal M,\bar a)\sim_\beta(\mathcal N,\bar b)\) implies
\((\mathcal M,\bar a)\sim_\alpha(\mathcal N,\bar b)\). On pointed structures
with tuples of a fixed length, each \(\sim_\alpha\) is reflexive, symmetric
and transitive, and it is preserved when the same permutation is applied to
both tuples.
\end{hblemma}

\begin{hbtheorem}{thm:p03-w03-p03-039-rank-characterization}{proved=annex:p03-w03-p03-039-rank-characterization}{Back-and-forth rank equals quantifier rank}
For all \(\tau\)-structures \(\mathcal M,\mathcal N\), tuples
\(\bar a\in M^n\), \(\bar b\in N^n\) and ordinals \(\alpha\),
\((\mathcal M,\bar a)\sim_\alpha(\mathcal N,\bar b)\) holds if and only if
\(\mathcal M\models\varphi(\bar a)\Leftrightarrow\mathcal N\models\varphi(\bar b)\)
for every \(\mathcal L_{\infty,\omega}\)-formula \(\varphi(\bar v)\), with
free variables among \(\bar v\), of quantifier rank at most \(\alpha\). In
particular \(\mathcal M\equiv_\alpha\mathcal N\) if and only if
\(\mathcal M\sim_\alpha\mathcal N\)
\cite[Thm.~2.13 and Ex.~2.14, p.~16]{Marker2014InfinitaryNotes}.
\end{hbtheorem}

The notes leave \cref{thm:p03-w03-p03-039-rank-characterization} as an
exercise; the annex proof uses the formulas of the next subsection for the
converse direction.

\subsection{Scott formulas, Scott rank and Scott's isomorphism theorem}
\label{sec:p03-w03-p03-039-scott}

\begin{hbdefinition}{def:p03-w03-p03-039-scott-formulas}{The formulas \(\Phi^{\mathcal M}_{\bar a,\alpha}\)}
For a \(\tau\)-structure \(\mathcal M\), a tuple \(\bar a\in M^{<\omega}\) and
an ordinal \(\alpha\), define \(\mathcal L_{\infty,\omega}\)-formulas
\(\Phi^{\mathcal M}_{\bar a,\alpha}(\bar v)\), with \(\lvert\bar v\rvert=\lvert\bar a\rvert\):
\[
 \Phi^{\mathcal M}_{\bar a,0}(\bar v)=\bigwedge_{\psi\in X}\psi(\bar v),
\]
where \(X\) is the set of atomic and negated atomic \(\tau\)-formulas
\(\psi(\bar v)\) with \(\mathcal M\models\psi(\bar a)\);
\(\Phi^{\mathcal M}_{\bar a,\alpha}=\bigwedge_{\beta<\alpha}\Phi^{\mathcal M}_{\bar a,\beta}\)
for a limit \(\alpha\); and
\[
 \Phi^{\mathcal M}_{\bar a,\beta+1}(\bar v)=
 \bigwedge_{b\in M}\exists w\,\Phi^{\mathcal M}_{\bar a,b,\beta}(\bar v,w)
 \;\land\;
 \forall w\bigvee_{b\in M}\Phi^{\mathcal M}_{\bar a,b,\beta}(\bar v,w),
\]
with \(w\) a variable not among \(\bar v\)
\cite[p.~16]{Marker2014InfinitaryNotes}. If \(\lvert M\rvert=\kappa\) is
infinite and \(\alpha<\kappa^+\), every conjunction and disjunction in
\(\Phi^{\mathcal M}_{\bar a,\alpha}\) has at most \(\kappa\) members, so
\(\Phi^{\mathcal M}_{\bar a,\alpha}\in\mathcal L_{\kappa^+,\omega}\)
\cite[p.~19]{Marker2014InfinitaryNotes}.
\end{hbdefinition}

\begin{hblemma}{thm:p03-w03-p03-039-scott-formula-lemma}{proved=annex:p03-w03-p03-039-scott-formula-lemma}{What the Scott formulas express}
For all \(\tau\)-structures \(\mathcal M,\mathcal N\), \(\bar a\in M^l\),
\(\bar b\in N^l\) and ordinals \(\alpha\),
\((\mathcal M,\bar a)\sim_\alpha(\mathcal N,\bar b)\) if and only if
\(\mathcal N\models\Phi^{\mathcal M}_{\bar a,\alpha}(\bar b)\). Moreover
\(\operatorname{qr}(\Phi^{\mathcal M}_{\bar a,\alpha})\le\alpha\)
\cite[Lemma~2.15, p.~17]{Marker2014InfinitaryNotes}.
\end{hblemma}

\begin{hblemma}{thm:p03-w03-p03-039-scott-rank}{proved=annex:p03-w03-p03-039-scott-rank}{The Scott rank of an infinite structure}
Let \(\mathcal M\) be an infinite \(\tau\)-structure. There is an ordinal
\(\alpha<\lvert M\rvert^+\) such that for all \(l\) and all
\(\bar a,\bar b\in M^l\), \((\mathcal M,\bar a)\sim_\alpha(\mathcal M,\bar b)\)
implies \((\mathcal M,\bar a)\sim_\beta(\mathcal M,\bar b)\) for every ordinal
\(\beta\). The least such \(\alpha\) is the \emph{Scott rank}
\(\operatorname{sr}(\mathcal M)\); it is also the least \(\alpha\) with
\(\Gamma_\alpha=\Gamma_{\alpha+1}\), where \(\Gamma_\alpha\) is the set of
pairs \((\bar a,\bar b)\) of equal-length tuples from \(M\) with
\((\mathcal M,\bar a)\sim_\alpha(\mathcal M,\bar b)\)
\cite[Lemma~2.17, pp.~17--18]{Marker2014InfinitaryNotes}.
\end{hblemma}

The notes assert \(\Gamma_\alpha\subseteq\Gamma_\beta\) for \(\alpha<\beta\)
and choose pairs in \(\Gamma_{\alpha+1}\setminus\Gamma_\alpha\)
\cite[p.~18]{Marker2014InfinitaryNotes}; by
\cref{thm:p03-w03-p03-039-approximant-basics} the inclusion runs the other
way, \(\Gamma_\beta\subseteq\Gamma_\alpha\), and the annex proof chooses
pairs in \(\Gamma_\alpha\setminus\Gamma_{\alpha+1}\).

\begin{hbdefinition}{def:p03-w03-p03-039-tuple-rank}{Ranks of tuples}
For a tuple \(\bar a\) from \(M\), \(r(\bar a)\) is the least ordinal
\(\alpha\) such that for all \(\bar b\) of the same length,
\((\mathcal M,\bar a)\sim_\alpha(\mathcal M,\bar b)\) implies
\((\mathcal M,\bar a)\sim_{\alpha+1}(\mathcal M,\bar b)\). The notes set
\(\operatorname{SR}(\mathcal M)=\sup\{r(\bar a)+1:\bar a\in M^{<\omega}\}\)
\cite[Def.~2.18, p.~18]{Marker2014InfinitaryNotes}.
\end{hbdefinition}

\begin{hblemma}{thm:p03-w03-p03-039-tuple-rank-bound}{proved=annex:p03-w03-p03-039-tuple-rank-bound}{Tuple ranks are bounded by the Scott rank}
For an infinite \(\mathcal M\), \(r(\bar a)\le\operatorname{sr}(\mathcal M)\)
for every tuple \(\bar a\), so
\(\sup_{\bar a}r(\bar a)\le\operatorname{sr}(\mathcal M)\).
\end{hblemma}

\begin{hbwarning}{def:p03-w03-p03-039-sr-equality}{The asserted equality \(\operatorname{sr}=\sup r\) is false}
The notes state, directly after Definition~2.18, that
\(\operatorname{sr}(\mathcal M)=\sup\{r(\bar a):\bar a\in M^{<\omega}\}\)
\cite[p.~18]{Marker2014InfinitaryNotes}, and give no argument. Only the
inequality of \cref{thm:p03-w03-p03-039-tuple-rank-bound} holds in general.
The next theorem gives a countable structure with
\(\sup_{\bar a}r(\bar a)=1<2=\operatorname{sr}(\mathcal M)\). The source of
the gap is that the one-step condition defining \(r(\bar a)\) is not
inherited upwards: a tuple can satisfy it at \(0\) and fail it at \(1\).
Nothing in this section depends on the equality; the Scott sentence below is
built at \(\operatorname{sr}(\mathcal M)\) itself.
\end{hbwarning}

\begin{hbtheorem}{thm:p03-w03-p03-039-sr-counterexample}{proved=annex:p03-w03-p03-039-sr-counterexample}{A structure whose Scott rank exceeds every tuple rank}
Let \(\tau=\{E,P\}\) with \(E\) binary and \(P\) unary. Let \(\mathcal M\)
have as universe the disjoint union of a five-element set
\(\{a_1,a_2,a_3,b_1,b_2\}\) and \(\omega\), with
\(P^{\mathcal M}=\omega\) and \(E^{\mathcal M}\) the equivalence relation on
\(\{a_1,a_2,a_3,b_1,b_2\}\) with the two classes \(\{a_1,a_2,a_3\}\) and
\(\{b_1,b_2\}\) (no pair involving a point of \(\omega\) is in
\(E^{\mathcal M}\)). Then \(r(\bar c)\le1\) for every tuple \(\bar c\) from
\(M\), \(r(a_1)=0\), and \(\operatorname{sr}(\mathcal M)=2\). Hence
\(\sup_{\bar c}r(\bar c)=1<\operatorname{sr}(\mathcal M)\), while
\(\operatorname{SR}(\mathcal M)=\sup_{\bar c}(r(\bar c)+1)=2\).
\end{hbtheorem}

For the infinite pure set of \cref{def:p03-w03-p03-039-pure-set} every tuple
has rank \(0\), so there \(\sup r=\operatorname{sr}=0\) and
\(\operatorname{SR}=1\): neither \(\sup r\) nor \(\operatorname{SR}\) equals
\(\operatorname{sr}\) for all structures.

\begin{hbdefinition}{def:p03-w03-p03-039-scott-sentence}{The Scott sentence}
Let \(\mathcal M\) be an infinite \(\tau\)-structure of cardinality \(\kappa\)
and let \(\alpha=\operatorname{sr}(\mathcal M)\). The \emph{Scott sentence} of
\(\mathcal M\) is
\[
 \Phi^{\mathcal M}=\Phi^{\mathcal M}_{\varnothing,\alpha}\land
 \bigwedge_{l=0}^{\infty}\ \bigwedge_{\bar a\in M^l}
 \forall\bar v\bigl(\Phi^{\mathcal M}_{\bar a,\alpha}(\bar v)\to
                    \Phi^{\mathcal M}_{\bar a,\alpha+1}(\bar v)\bigr).
\]
Since \(\alpha<\kappa^+\), \(\Phi^{\mathcal M}\in\mathcal L_{\kappa^+,\omega}\);
if \(\tau\) and \(\mathcal M\) are countable, \(\Phi^{\mathcal M}\in\mathcal L_{\omega_1,\omega}\)
\cite[p.~19]{Marker2014InfinitaryNotes}.
\end{hbdefinition}

\begin{hbtheorem}{thm:p03-w03-p03-039-scott}{proved=annex:p03-w03-p03-039-scott}{Scott's isomorphism theorem}
Let \(\mathcal M\) be an infinite \(\tau\)-structure and \(\mathcal N\) a
\(\tau\)-structure. Then \(\mathcal N\equiv_{\infty,\omega}\mathcal M\) if and
only if \(\mathcal N\models\Phi^{\mathcal M}\). In particular, if \(\tau\),
\(\mathcal M\) and \(\mathcal N\) are countable, then
\(\Phi^{\mathcal M}\in\mathcal L_{\omega_1,\omega}\) and
\(\mathcal N\models\Phi^{\mathcal M}\) implies \(\mathcal N\cong\mathcal M\)
\cite[Thm.~2.22, pp.~19--20]{Marker2014InfinitaryNotes}.
\end{hbtheorem}

\begin{hbdefinition}{def:p03-w03-p03-039-complete}{Complete sentences}
A satisfiable sentence \(\varphi\in\mathcal L_{\infty,\omega}\) is
\emph{complete} if for every sentence \(\psi\in\mathcal L_{\infty,\omega}\),
\(\varphi\models\psi\) or \(\varphi\models\neg\psi\)
\cite[Def.~2.25, p.~20]{Marker2014InfinitaryNotes}.
\end{hbdefinition}

\begin{hbcorollary}{thm:p03-w03-p03-039-scott-complete}{proved=annex:p03-w03-p03-039-scott-complete}{Scott sentences are complete}
For every infinite \(\tau\)-structure \(\mathcal M\), the Scott sentence
\(\Phi^{\mathcal M}\) is complete
\cite[Cor.~2.26, p.~20]{Marker2014InfinitaryNotes}.
\end{hbcorollary}

\begin{hbexample}{def:p03-w03-p03-039-pure-set}{The Scott sentence of an infinite pure set}
Let \(\tau=\varnothing\) and let \(\mathcal M\) be an infinite set. Then
\((\mathcal M,\bar a)\sim_0(\mathcal M,\bar b)\) says that \(\bar a\) and
\(\bar b\) have the same equality pattern: \(a_i=a_j\) exactly when
\(b_i=b_j\). This already implies \(\sim_1\): given \(c\in M\), put
\(d=b_i\) if \(c=a_i\), and otherwise let \(d\) be any element outside the
finite set \(\{b_1,\ldots,b_n\}\), which exists because \(M\) is infinite;
then \(\bar ac\) and \(\bar bd\) have the same equality pattern, and the
back clause is symmetric. Hence \(\Gamma_0=\Gamma_1\) and
\(\operatorname{sr}(\mathcal M)=0\). The conjunct
\(\Phi^{\mathcal M}_{\varnothing,0}\) is the empty conjunction, true in every
structure, and the conjunct for an injective tuple \(\bar a\) of length
\(l\) says that whenever \(\bar v\) is injective there is a \(w\) outside
\(\bar v\). A structure \(\mathcal N\) satisfies all these conjuncts exactly
when no injective \(l\)-tuple exhausts \(N\) for any \(l\), that is, when
\(N\) is infinite. So the countable models of \(\Phi^{\mathcal M}\) are the
countably infinite sets, as \cref{thm:p03-w03-p03-039-scott} predicts, and
the uncountable models are the uncountable sets, all
\(\equiv_{\infty,\omega}\)-equivalent to \(\mathcal M\).
\end{hbexample}

\subsection{Consistency properties and model existence}
\label{sec:p03-w03-p03-039-model-existence}

Throughout this subsection \(\tau\) is countable and \(C\) is an infinite
set of constant symbols of \(\tau\), not necessarily all of them; since
\(\tau\) is countable, \(C\) is countably infinite. The notion of a
consistency property is attributed in the notes to Makkai
\cite[p.~39]{Marker2014InfinitaryNotes}. The notes also state that no useful
completeness theorem is available in this setting
\cite[p.~39]{Marker2014InfinitaryNotes}; the model existence theorem is the
tool used in its place, and no infinitary proof calculus is part of this
presentation.

\begin{hbdefinition}{def:p03-w03-p03-039-consistency-property}{Consistency property}
A \emph{consistency property} is a collection \(\Sigma\) of countable sets of
\(\mathcal L_{\omega_1,\omega}(\tau)\)-sentences such that, for every
\(\sigma\in\Sigma\):
\begin{enumerate}
\item[C0] if \(\mu\subseteq\sigma\) then \(\mu\in\Sigma\);
\item[C1] if \(\varphi\in\sigma\) then \(\neg\varphi\notin\sigma\);
\item[C2] if \(\neg\varphi\in\sigma\) there is \(\mu\in\Sigma\) with
      \(\sigma\cup\{{\sim}\varphi\}\subseteq\mu\);
\item[C3] if \(\bigwedge_{\varphi\in X}\varphi\in\sigma\), then for every
      \(\varphi\in X\) there is \(\mu\in\Sigma\) with
      \(\sigma\cup\{\varphi\}\subseteq\mu\);
\item[C4] if \(\bigvee_{\varphi\in X}\varphi\in\sigma\), there are
      \(\mu\in\Sigma\) and \(\varphi\in X\) with \(\sigma\cup\{\varphi\}\subseteq\mu\);
\item[C5] if \(\forall v\,\varphi(v)\in\sigma\), then for every \(c\in C\)
      there is \(\mu\in\Sigma\) with \(\sigma\cup\{\varphi(c)\}\subseteq\mu\);
\item[C6] if \(\exists v\,\varphi(v)\in\sigma\), there are \(\mu\in\Sigma\)
      and \(c\in C\) with \(\sigma\cup\{\varphi(c)\}\subseteq\mu\);
\item[C7] for every closed term \(t\) and all \(c,d\in C\):
      (a) if \(c=d\in\sigma\), there is \(\mu\in\Sigma\) with
      \(\sigma\cup\{d=c\}\subseteq\mu\);
      (b) if \(c=t\in\sigma\) and \(\varphi(t)\in\sigma\), there is
      \(\mu\in\Sigma\) with \(\sigma\cup\{\varphi(c)\}\subseteq\mu\);
      (c) there are \(\mu\in\Sigma\) and \(e\in C\) with
      \(\sigma\cup\{e=t\}\subseteq\mu\)
\end{enumerate}
\cite[Def.~4.1, p.~39]{Marker2014InfinitaryNotes}.
\end{hbdefinition}

\begin{hblemma}{thm:p03-w03-p03-039-consistency-lemma}{proved=annex:p03-w03-p03-039-consistency-lemma}{Equality and modus ponens inside a consistency property}
Let \(\Sigma\) be a consistency property, \(\sigma\in\Sigma\) and
\(c,d,e\in C\).
\begin{enumerate}
\item There is \(\mu\in\Sigma\) with \(\sigma\subseteq\mu\) and \(c=c\in\mu\).
\item If \(c=d\in\sigma\) and \(d=e\in\sigma\), then \(\sigma\cup\{c=e\}\in\Sigma\).
\item If \(\varphi\in\sigma\) and \(\varphi\to\psi\in\sigma\), there is
      \(\mu\in\Sigma\) with \(\sigma\subseteq\mu\) and \(\psi\in\mu\).
\end{enumerate}
\cite[Lemma~4.2 and Ex.~4.3, pp.~39--40]{Marker2014InfinitaryNotes}.
\end{hblemma}

\begin{hbtheorem}{thm:p03-w03-p03-039-model-existence}{proved=annex:p03-w03-p03-039-model-existence}{Model existence theorem}
If \(\Sigma\) is a consistency property and \(\sigma\in\Sigma\), there is a
countable \(\tau\)-structure \(\mathcal M\) with \(\mathcal M\models\sigma\)
\cite[Thm.~4.6, pp.~40--43]{Marker2014InfinitaryNotes}.
\end{hbtheorem}

\begin{hbtheorem}{thm:p03-w03-p03-039-extended-model-existence}{proved=annex:p03-w03-p03-039-extended-model-existence}{Extended model existence theorem}
Let \(T\) be a countable set of \(\mathcal L_{\omega_1,\omega}(\tau)\)-sentences
and let \(\Sigma\) be a \emph{nonempty} consistency property such that
\(\sigma\cup\{\varphi\}\in\Sigma\) for all \(\sigma\in\Sigma\) and
\(\varphi\in T\). Then \(T\) has a countable model
\cite[Thm.~4.8, p.~43]{Marker2014InfinitaryNotes}.
\end{hbtheorem}

\begin{hbwarning}{def:p03-w03-p03-039-nonempty}{Where the Henkin argument needs care}
The notes state \cref{thm:p03-w03-p03-039-extended-model-existence}
without the word ``nonempty''. The empty collection satisfies C0--C7 and
the closure hypothesis vacuously, and \(T=\{\neg\,\forall v\,v=v\}\) has no
model, so nonemptiness is a needed hypothesis. Two further points of the
printed proof of \cref{thm:p03-w03-p03-039-model-existence} are made
explicit in the annex: the set \(\Lambda\) of sentences that the
construction may add must also contain the atomic and negated atomic
sentences, because the equations \(c=t\) produced by C7(c) and the
substitution instances produced by C7(b) are used later; and the induction
proving that the canonical structure satisfies the constructed set is on a
syntactic depth for which \({\sim}\theta\) is below \(\neg\psi\) whenever
\(\theta\) is an immediate constituent of \(\psi\), not on \({\sim}\psi\)
itself.
\end{hbwarning}

\subsection{Definability of the approximants, free abelian groups and complete sentences}
\label{sec:p03-w03-p03-039-complete-sentences}

Throughout this subsection \(\tau\) is countable unless a statement says
otherwise. A sentence is \emph{\(\aleph_0\)-categorical} if it has, up to
isomorphism, exactly one model of cardinality at most \(\aleph_0\), and
\emph{\(\kappa\)-categorical} if it has exactly one model of cardinality
\(\kappa\).

\begin{hblemma}{thm:p03-w03-p03-039-countable-fragment-size}{proved=annex:p03-w03-p03-039-countable-fragment-size}{Fragments of a given size}
Let \(\kappa\) be an infinite cardinal, \(\lvert\tau\rvert\le\kappa\), \(T\) a
set of at most \(\kappa\) formulas of \(\mathcal L_{\kappa^+,\omega}(\tau)\)
each with finitely many free variables, and \(V\) a set of \(\kappa\)
variables containing every variable occurring in \(T\). There is a fragment
\(A\supseteq T\) with witness \(V\) and \(\lvert A\rvert\le\kappa\) all of
whose formulas have finitely many free variables. For \(\kappa=\aleph_0\) and
\(T\) a set of sentences, the least fragment of
\cref{thm:p03-w03-p03-039-countable-fragment} is such a fragment.
\end{hblemma}

This is the form of \cite[Exs.~1.19 and~1.20, p.~9]{Marker2014InfinitaryNotes}
used below; the hypothesis \(\lvert\tau\rvert\le\kappa\) is the one
\cref{def:p03-w03-p03-039-fragment-size} shows to be needed.

\begin{hbtheorem}{thm:p03-w03-p03-039-definable-approximants}{proved=annex:p03-w03-p03-039-definable-approximants}{The approximants and the Scott rank are definable}
\begin{enumerate}
\item For all \(\beta<\omega_1\) and \(n\in\omega\) there is an
      \(\mathcal L_{\omega_1,\omega}(\tau)\)-formula \(S^n_\beta(\bar v,\bar w)\),
      \(\lvert\bar v\rvert=\lvert\bar w\rvert=n\), such that for every
      \(\tau\)-structure \(\mathcal M\) and all \(\bar a,\bar b\in M^n\),
      \(\mathcal M\models S^n_\beta(\bar a,\bar b)\) if and only if
      \((\mathcal M,\bar a)\sim_\beta(\mathcal M,\bar b)\).
\item For every \(\tau\)-structure \(\mathcal M\) and tuple \(\bar a\) the rank
      \(r(\bar a)\) of \cref{def:p03-w03-p03-039-tuple-rank} exists, and for
      all \(\alpha<\omega_1\) there is \(\sigma_\alpha\in\mathcal L_{\omega_1,\omega}(\tau)\)
      such that \(\mathcal M\models\sigma_\alpha\) if and only if
      \(\operatorname{SR}(\mathcal M)\ge\alpha\).
\end{enumerate}
\cite[Ex.~2.21, p.~19]{Marker2014InfinitaryNotes}.
\end{hbtheorem}

The notes print the right-hand side of item~1 as
\((\mathcal M,\bar a)\sim_\alpha(\mathcal N,\bar b)\)
\cite[Ex.~2.21(a), p.~19]{Marker2014InfinitaryNotes}; both tuples are from
the one structure \(\mathcal M\) and the index is \(\beta\), as above. Item~2
uses \(\operatorname{SR}(\mathcal M)=\sup\{r(\bar a)+1\}\) of
\cref{def:p03-w03-p03-039-tuple-rank}.

\begin{hblemma}{thm:p03-w03-p03-039-countable-substructures}{proved=annex:p03-w03-p03-039-countable-substructures}{Countably generated substructures transfer}
If \(\mathcal M\equiv_{\infty,\omega}\mathcal N\) (for any vocabulary), then
every countably generated substructure of \(\mathcal M\) is isomorphic to a
substructure of \(\mathcal N\)
\cite[Ex.~2.5, p.~14]{Marker2014InfinitaryNotes}.
\end{hblemma}

The notes print ``\(\mathcal M_0\subseteq\mathcal N\)'' for the given
substructure \cite[Ex.~2.5, p.~14]{Marker2014InfinitaryNotes}; it is a
substructure of \(\mathcal M\).

\begin{hbdefinition}{def:p03-w03-p03-039-aleph-one-free}{\(\aleph_1\)-free abelian groups}
An abelian group, in the vocabulary \(\{+,-,0\}\), is \emph{\(\aleph_1\)-free}
if every countable subgroup of it is free abelian
\cite[Ex.~2.6, p.~14]{Marker2014InfinitaryNotes}. A free abelian group is
one with a basis, possibly empty.
\end{hbdefinition}

\begin{hblemma}{thm:p03-w03-p03-039-subgroups-free}{proved=annex:p03-w03-p03-039-subgroups-free}{Countable subgroups of free abelian groups are free}
Every countable subgroup of a free abelian group is free abelian. In
particular, if \(G\) is free abelian and \(H\equiv_{\infty,\omega}G\), then
\(H\) is \(\aleph_1\)-free
\cite[Ex.~2.6, p.~14]{Marker2014InfinitaryNotes}.
\end{hblemma}

\begin{hbcorollary}{thm:p03-w03-p03-039-aleph-one-free}{proved=annex:p03-w03-p03-039-aleph-one-free}{\(\aleph_1\)-free groups and the free group of countable rank}
Let \(G\) be the free abelian group on countably infinitely many generators
and let \(H\) be an \(\aleph_1\)-free abelian group that is not free on a
finite set of generators. Then \(G\equiv_{\infty,\omega}H\)
\cite[Cor.~2.23, p.~20]{Marker2014InfinitaryNotes}.
\end{hbcorollary}

The notes prove \cref{thm:p03-w03-p03-039-aleph-one-free} from Scott's
theorem and the downward L\"owenheim--Skolem theorem; the annex follows that
route and writes out the sentence saying that a group is not free on finitely
many generators, which the printed proof uses without displaying.

\begin{hbtheorem}{thm:p03-w03-p03-039-aleph-one-free-axiom}{proved=annex:p03-w03-p03-039-aleph-one-free-axiom}{\(\aleph_1\)-freeness is axiomatizable}
The class of \(\aleph_1\)-free abelian groups is the class of models of one
sentence of \(\mathcal L_{\omega_1,\omega}(\{+,-,0\})\)
\cite[Ex.~2.24, p.~20]{Marker2014InfinitaryNotes}.
\end{hbtheorem}

\begin{hblemma}{thm:p03-w03-p03-039-finite-structures}{proved=annex:p03-w03-p03-039-finite-structures}{Finite structures and countable models}
\begin{enumerate}
\item Every finite \(\tau\)-structure \(\mathcal M\) has a sentence
      \(\chi_{\mathcal M}\in\mathcal L_{\omega_1,\omega}(\tau)\) whose models are
      exactly the structures isomorphic to \(\mathcal M\).
\item Every satisfiable sentence of \(\mathcal L_{\omega_1,\omega}(\tau)\) has
      a model of cardinality at most \(\aleph_0\).
\end{enumerate}
\end{hblemma}

\begin{hbtheorem}{thm:p03-w03-p03-039-complete-extension}{proved=annex:p03-w03-p03-039-complete-extension}{Every satisfiable sentence has a complete strengthening}
For every satisfiable \(\psi\in\mathcal L_{\omega_1,\omega}(\tau)\) there is a
complete \(\varphi\in\mathcal L_{\omega_1,\omega}(\tau)\) with
\(\varphi\models\psi\)
\cite[Ex.~2.27, p.~20]{Marker2014InfinitaryNotes}.
\end{hbtheorem}

\begin{hbwarning}{def:p03-w03-p03-039-countable-vocabulary}{The countability of \(\tau\) is needed}
\Cref{thm:p03-w03-p03-039-complete-extension} fails for uncountable
\(\tau\). Let \(\tau=\{P_\xi:\xi<\omega_1\}\) with each \(P_\xi\) unary and let
\(\varphi\in\mathcal L_{\omega_1,\omega}(\tau)\) be satisfiable. Only
countably many symbols occur in \(\varphi\); pick \(P_\xi\) not among them
and a model \(\mathcal M\) of \(\varphi\). Changing the interpretation of
\(P_\xi\) to \(\varnothing\), and then to \(M\), does not affect the truth of
\(\varphi\), so \(\varphi\) has a model of \(\exists x\,P_\xi(x)\) and a model of
its negation, and is not complete.
\end{hbwarning}

\begin{hbtheorem}{thm:p03-w03-p03-039-complete-categorical}{proved=annex:p03-w03-p03-039-complete-categorical}{Complete means \(\aleph_0\)-categorical}
A satisfiable \(\varphi\in\mathcal L_{\omega_1,\omega}(\tau)\) is complete if
and only if it is \(\aleph_0\)-categorical
\cite[Ex.~2.28, p.~20]{Marker2014InfinitaryNotes}.
\end{hbtheorem}

\begin{hbtheorem}{thm:p03-w03-p03-039-categorical-equivalence}{proved=annex:p03-w03-p03-039-categorical-equivalence}{Above a categoricity cardinal all models are \(\equiv_{\infty,\omega}\)}
Let \(\psi\in\mathcal L_{\omega_1,\omega}(\tau)\) be \(\kappa\)-categorical for
an infinite cardinal \(\kappa\). If \(\mathcal M,\mathcal N\models\psi\) and
\(\lvert M\rvert,\lvert N\rvert\ge\kappa\), then
\(\mathcal M\equiv_{\infty,\omega}\mathcal N\)
\cite[Ex.~2.29, p.~21]{Marker2014InfinitaryNotes}.
\end{hbtheorem}

The proof applies the downward L\"owenheim--Skolem theorem to a fragment of
cardinality \(\kappa\) containing \(\psi\) and the Scott sentence of the model
of cardinality \(\kappa\); the notes compare the statement with a later
exercise on Hanf numbers \cite[Ex.~2.29, p.~21]{Marker2014InfinitaryNotes}.

\begin{hbexample}{def:p03-w03-p03-039-categorical-incomplete}{An \(\aleph_1\)-categorical sentence that is not complete}
In the vocabulary \(\{<\}\) let
\(\varphi=\sigma_{\mathbb Z}\vee\forall x\forall y\,\neg\,x<y\), with
\(\sigma_{\mathbb Z}\) the sentence of
\cref{thm:p03-w03-p03-039-axiomatizable}~iv). Every model of
\(\sigma_{\mathbb Z}\) is countable
(\cref{thm:p03-w03-p03-039-upward-ls-failure}), so a model of cardinality
\(\aleph_1\) satisfies the second disjunct: it is a set of \(\aleph_1\)
elements with the empty relation, and any two such are isomorphic by any
bijection. So \(\varphi\) is \(\aleph_1\)-categorical. It is not complete:
\((\mathbb Z,<)\models\varphi\wedge\sigma_{\mathbb Z}\) and the set \(\omega\)
with the empty relation satisfies \(\varphi\wedge\neg\sigma_{\mathbb Z}\). This
answers \cite[Ex.~2.30, p.~21]{Marker2014InfinitaryNotes}, and shows that
\cref{thm:p03-w03-p03-039-categorical-equivalence} says nothing below
\(\kappa\).
\end{hbexample}

\subsection{Uncountable models of complete sentences and countable approximations}
\label{sec:p03-w03-p03-039-approximations}

This subsection treats the end of the notes' second chapter: when a complete
sentence has an uncountable model, the reduction of complete sentences to
atomic models of first-order theories, and Kueker's description of
\(\equiv_{\infty,\omega}\) through countable approximations
\cite[pp.~21--26]{Marker2014InfinitaryNotes}. For tuples
\(\bar a\in M^n\), \(\bar b\in N^n\) write
\((\mathcal M,\bar a)\equiv_{\infty,\omega}(\mathcal N,\bar b)\) when
\(\mathcal M\models\varphi(\bar a)\Leftrightarrow\mathcal N\models\varphi(\bar b)\)
for every \(\mathcal L_{\infty,\omega}\)-formula \(\varphi(\bar v)\) with free
variables among \(\bar v\). For a fragment \(A\), a map
\(f:M\to N\) is \emph{\(A\)-elementary} if
\(\mathcal M\models\varphi(\bar a)\Leftrightarrow\mathcal N\models\varphi(f(\bar a))\)
for all \(\varphi(\bar v)\in A\) and \(\bar a\) from \(M\); it is
\emph{\(\mathcal L_{\omega_1,\omega}\)-elementary} if this holds for every
\(\mathcal L_{\omega_1,\omega}\)-formula with finitely many free variables.
An \(A\)-elementary map is injective and is an embedding, since \(A\)
contains the atomic formulas.

\begin{hblemma}{thm:p03-w03-p03-039-elementary-chain}{proved=annex:p03-w03-p03-039-elementary-chain}{Unions of \(A\)-elementary chains}
Let \(A\) be a fragment all of whose formulas have finitely many free
variables, \((I,<)\) a nonempty linear order and \((\mathcal M_i:i\in I)\)
\(\tau\)-structures with \(\mathcal M_i\prec_A\mathcal M_j\) for \(i<j\). Let
\(\mathcal M\) be the union of the chain. Then
\(\mathcal M_i\prec_A\mathcal M\) for every \(i\in I\)
\cite[Ex.~1.28, p.~11]{Marker2014InfinitaryNotes}.
\end{hblemma}

\begin{hbtheorem}{thm:p03-w03-p03-039-uncountable-models}{proved=annex:p03-w03-p03-039-uncountable-models}{When a complete sentence has an uncountable model}
Let \(\tau\) be countable, \(\sigma\in\mathcal L_{\omega_1,\omega}(\tau)\) a
complete sentence, \(\mathcal M\) its model of cardinality at most
\(\aleph_0\) (unique up to isomorphism by
\cref{thm:p03-w03-p03-039-complete-categorical}), and \(A\) a countable
fragment of \(\mathcal L_{\omega_1,\omega}(\tau)\) with \(\sigma\in A\), all
of whose formulas have finitely many free variables. The following are
equivalent:
\begin{enumerate}
\item \(\sigma\) has an uncountable model;
\item there is a non-surjective \(\mathcal L_{\omega_1,\omega}\)-elementary
      map of \(\mathcal M\) into itself;
\item there is a non-surjective \(A\)-elementary map of \(\mathcal M\) into
      itself
\end{enumerate}
\cite[Ex.~2.31, p.~21]{Marker2014InfinitaryNotes}.
\end{hbtheorem}

The notes write ``a countable fragment containing \(\varphi\)'' in item~3;
the sentence meant is \(\sigma\). The proof of 1\(\Rightarrow\)2 in fact
produces a map preserving every \(\mathcal L_{\infty,\omega}\)-formula. The
notes add, without proof, that Gao has shown these conditions equivalent to
the nonexistence of a complete left-invariant metric on the automorphism
group of \(\mathcal M\) \cite[p.~21]{Marker2014InfinitaryNotes}; that
equivalence is not proved here.

\begin{hbtheorem}{thm:p03-w03-p03-039-complete-atomic}{proved=annex:p03-w03-p03-039-complete-atomic}{Complete sentences and atomic models}
Let \(\tau\) be countable and \(\varphi\in\mathcal L_{\omega_1,\omega}(\tau)\)
a complete sentence. There are a countable vocabulary \(\tau^*\supseteq\tau\)
and a complete first-order \(\tau^*\)-theory \(T^*\) such that a
\(\tau\)-structure \(\mathcal N\) is a model of \(\varphi\) if and only if
\(\mathcal N\) is the \(\tau\)-reduct of an atomic model of \(T^*\)
\cite[Thm.~2.32, p.~21]{Marker2014InfinitaryNotes}.
\end{hbtheorem}

The proof refines Chang's reduction (\cref{thm:p03-w03-p03-039-chang}): it
takes the Chang expansion for a countable fragment containing \(\varphi\)
and the Scott formulas \(\Phi^{\mathcal M}_{\bar b,\alpha}\) of the countable
model \(\mathcal M\) at \(\alpha=\operatorname{sr}(\mathcal M)\), and lets
\(T^*\) be the complete theory of the natural expansion of \(\mathcal M\).
The family of omitted types disappears from the statement: an atomic model
omits it automatically. The notes' proof treats an infinite \(\mathcal M\);
the annex adds the case of a finite one, where the Scott rank is not
defined.

\begin{hbdefinition}{def:p03-w03-p03-039-a-types}{\(A\)-types}
Let \(A\) be a fragment of \(\mathcal L_{\omega_1,\omega}(\tau)\). For a
\(\tau\)-structure \(\mathcal M\) and \(\bar a=(a_1,\ldots,a_n)\) from \(M\),
the \emph{\(A\)-type} of \(\bar a\) is
\(\operatorname{tp}^{\mathcal M}(\bar a,A)=\{\varphi(\bar v)\in A:\mathcal M\models\varphi(\bar a)\}\),
and \(S_n(A,\varphi)\) is the set of \(A\)-types of \(n\)-tuples realized in
some model of the sentence \(\varphi\)
\cite[Def.~2.33, p.~22]{Marker2014InfinitaryNotes}. The
\(\mathcal L_{\omega_1,\omega}\)-type of \(\bar a\) is the set of
\(\mathcal L_{\omega_1,\omega}(\tau)\)-formulas with free variables among
\(\bar v\) that \(\bar a\) satisfies.
\end{hbdefinition}

The notes print ``realized in some model of \(T\)'' in the definition of
\(S_n(A,\varphi)\); the only theory in sight is the sentence \(\varphi\).

\begin{hbtheorem}{thm:p03-w03-p03-039-countably-many-types}{proved=annex:p03-w03-p03-039-countably-many-types}{Countably many types and countable \(\equiv_{\infty,\omega}\)-representatives}
Let \(\tau\) be countable and \(\mathcal M\) a \(\tau\)-structure of any
cardinality.
\begin{enumerate}
\item If \(\mathcal M\equiv_{\infty,\omega}\mathcal N\) for some countable
      \(\mathcal N\), then \(\mathcal M\) realizes only countably many
      \(\mathcal L_{\omega_1,\omega}\)-types.
\item If \(\mathcal M\) is infinite and realizes only countably many
      \(\mathcal L_{\omega_1,\omega}\)-types, then
      \(\operatorname{sr}(\mathcal M)<\omega_1\).
\item If \(\mathcal M\) realizes only countably many
      \(\mathcal L_{\omega_1,\omega}\)-types, then there is a countable
      \(\mathcal N\equiv_{\infty,\omega}\mathcal M\).
\end{enumerate}
\cite[Ex.~2.34, p.~22]{Marker2014InfinitaryNotes}.
\end{hbtheorem}

The notes state the exercise for a countable \(\mathcal M\). With that
hypothesis all three items are immediate: \(\mathcal M\) has countably many
tuples, \(\operatorname{sr}(\mathcal M)<\lvert M\rvert^+=\omega_1\) by
\cref{thm:p03-w03-p03-039-scott-rank}, and \(\mathcal N=\mathcal M\) serves.
The statement above, for a structure of any cardinality, is the one with
content, and it is the one proved in the annex. Items~1 and~3 together say
that \(\mathcal M\) is \(\equiv_{\infty,\omega}\)-equivalent to a countable
structure exactly when it realizes countably many
\(\mathcal L_{\omega_1,\omega}\)-types.

\paragraph{Countable approximations.} The rest of the subsection presents
Kueker's characterization of \(\equiv_{\infty,\omega}\)
\cite[Sec.~2.3, pp.~22--26]{Marker2014InfinitaryNotes}. The background
theory is ZFC; the axiom of choice is used to pick members of closed
unbounded families and to know that countable unions of countable sets are
countable. For a set \(X\), \(\mathcal P_{\omega_1}(X)\) is the set of
countable (finite or countably infinite) subsets of \(X\).

\begin{hbdefinition}{def:p03-w03-p03-039-club}{Closed unbounded families and the filter \(D(X)\)}
A family \(\mathcal A\subseteq\mathcal P_{\omega_1}(X)\) is \emph{closed} if
\(\bigcup_{n}a_n\in\mathcal A\) whenever \(a_0\subseteq a_1\subseteq\cdots\)
are members of \(\mathcal A\), and \emph{unbounded} if for every
\(a\in\mathcal P_{\omega_1}(X)\) there is \(b\in\mathcal A\) with
\(a\subseteq b\); a \emph{club} is a closed unbounded family.
\(D(X)\) is the set of \(\mathcal A\subseteq\mathcal P_{\omega_1}(X)\) that
include a club. A property of \(a\in\mathcal P_{\omega_1}(X)\) holds
\emph{almost everywhere} (a.e.) if the set of \(a\) with the property is in
\(D(X)\) \cite[Def.~2.35, p.~22; p.~23]{Marker2014InfinitaryNotes}.
\end{hbdefinition}

The notes print the closure condition as ``\(\bigcup a_n\in X\) for any
\(a_1\subset a_2\subset\cdots\) in \(X\)''; the family meant is
\(\mathcal A\). Whether the chains are strictly increasing makes no
difference: a non-strict chain is eventually constant, with its union a
member, or has a strictly increasing subsequence with the same union. In the
proof of Lemma~2.36 the notes write the new set as an intersection
\(\bigcap a_n\); it is the union.

\begin{hblemma}{thm:p03-w03-p03-039-club-filter}{proved=annex:p03-w03-p03-039-club-filter}{Clubs are closed under countable and diagonal intersections}
Let \(X\) be a set.
\begin{enumerate}
\item The intersection of countably many clubs is a club, and \(D(X)\) is a
      countably complete proper filter on \(\mathcal P_{\omega_1}(X)\): it
      contains \(\mathcal P_{\omega_1}(X)\), does not contain \(\varnothing\),
      and is closed upwards and under countable intersections.
\item For \(I\subseteq X\) and clubs \(\mathcal A_i\) \((i\in I)\), the
      diagonal intersection
      \(\triangle_{i\in I}\mathcal A_i=\{a\in\mathcal P_{\omega_1}(X):a\in\mathcal A_i\text{ for all }i\in a\cap I\}\)
      is a club.
\item For every countable \(s\subseteq X\), \(\{a:s\subseteq a\}\) is a club.
\end{enumerate}
\cite[Lemmas~2.36 and~2.38, Ex.~2.37, pp.~22--23]{Marker2014InfinitaryNotes}.
\end{hblemma}

\begin{hbtheorem}{thm:p03-w03-p03-039-kueker-game}{proved=annex:p03-w03-p03-039-kueker-game}{Kueker's game}
Let \(X\) be nonempty and \(\mathcal A\subseteq\mathcal P_{\omega_1}(X)\). In
the game \(G(\mathcal A)\) Players~I and~II alternately play elements of
\(X\), Player~I playing \(c_0,c_1,\ldots\) and Player~II answering
\(d_0,d_1,\ldots\); Player~II wins the play if
\(\{c_i,d_i:i\in\omega\}\in\mathcal A\). Player~II has a winning strategy
in \(G(\mathcal A)\) if and only if \(\mathcal A\in D(X)\)
\cite[Ex.~2.39, p.~23]{Marker2014InfinitaryNotes}.
\end{hbtheorem}

Let \(\tau\) be countable and \(\mathcal M\) a \(\tau\)-structure. A set \(X\)
is \emph{large enough to approximate} \(\mathcal M\) if \(M\subseteq X\); for
\(a\in\mathcal P_{\omega_1}(X)\) write \(\mathcal M^a\) for \(M\cap a\) with
the structure induced from \(\mathcal M\) whenever this set is the universe
of a substructure \cite[pp.~23--24]{Marker2014InfinitaryNotes}.

\begin{hblemma}{thm:p03-w03-p03-039-substructure-ae}{proved=annex:p03-w03-p03-039-substructure-ae}{Approximations are substructures almost everywhere}
If \(\tau\) is countable and \(X\) is large enough to approximate
\(\mathcal M\), then the set of \(a\in\mathcal P_{\omega_1}(X)\) such that
\(M\cap a\) is the universe of a substructure of \(\mathcal M\) is in
\(D(X)\) \cite[Ex.~2.40, p.~24]{Marker2014InfinitaryNotes}.
\end{hblemma}

\begin{hbdefinition}{def:p03-w03-p03-039-approximation}{The approximation \(\varphi^a\)}
For approximations, a conjunction or disjunction is read as an indexed family
\(\bigwedge_{i\in I}\varphi_i\), \(\bigvee_{i\in I}\varphi_i\), with \(I\) a
set; the conjunction \(\bigwedge_{\varphi\in Y}\varphi\) of
\cref{def:p03-w03-p03-039-language} is the family indexed by \(Y\) itself,
and satisfaction does not depend on the indexing. For a formula \(\varphi\)
of \(\mathcal L_{\infty,\omega}\) and \(a\in\mathcal P_{\omega_1}(X)\) the
approximation \(\varphi^a\) is defined by recursion: \(\varphi^a=\varphi\)
for atomic \(\varphi\); \((\neg\varphi)^a=\neg\varphi^a\);
\((\forall v\,\varphi)^a=\forall v\,\varphi^a\) and
\((\exists v\,\varphi)^a=\exists v\,\varphi^a\);
\((\bigwedge_{i\in I}\varphi_i)^a=\bigwedge_{i\in I\cap a}\varphi_i^a\) and
\((\bigvee_{i\in I}\varphi_i)^a=\bigvee_{i\in I\cap a}\varphi_i^a\).
\(X\) is \emph{large enough to approximate} \(\varphi\) if \(I\subseteq X\)
for every conjunction or disjunction \(\bigwedge_{i\in I}\varphi_i\),
\(\bigvee_{i\in I}\varphi_i\) that is a subformula of \(\varphi\). Each
\(\varphi^a\) is an \(\mathcal L_{\omega_1,\omega}\)-formula, since \(a\) is
countable \cite[Def.~2.41, p.~24]{Marker2014InfinitaryNotes}.
\end{hbdefinition}

The notes require \(I\subseteq X\) for every conjunction and disjunction in
the least fragment containing \(\varphi\); that requirement includes the one
above, since fragments are closed under subformulas, and only subformulas
enter the proofs. The negation normal form of the notes' Ex.~1.21 is
obtained from the formal negation \({\sim}\) of
\cref{def:p03-w03-p03-039-formal-negation}, read with the same index sets
(\({\sim}\bigwedge_{i\in I}\varphi_i=\bigvee_{i\in I}{\sim}\varphi_i\) and
dually): \(\operatorname{nnf}\) is the identity on atomic formulas,
\(\operatorname{nnf}(\neg\varphi)={\sim}\operatorname{nnf}(\varphi)\), and
\(\operatorname{nnf}\) commutes with \(\bigwedge_{i\in I}\),
\(\bigvee_{i\in I}\), \(\exists v\) and \(\forall v\).

\begin{hblemma}{thm:p03-w03-p03-039-approximation-nnf}{proved=annex:p03-w03-p03-039-approximation-nnf}{Approximation commutes with negation normal form}
For every formula \(\varphi\) of \(\mathcal L_{\infty,\omega}\), the formula
\(\operatorname{nnf}(\varphi)\) applies \(\neg\) only to atomic formulas, has
the same free variables and the same index sets as \(\varphi\), and
\(\models\varphi\leftrightarrow\operatorname{nnf}(\varphi)\). For every
\(a\), \(\operatorname{nnf}(\varphi)^a=\operatorname{nnf}(\varphi^a)\), so
\(\models\varphi^a\leftrightarrow\operatorname{nnf}(\varphi)^a\) for every
\(a\) \cite[Exs.~1.21 and~2.42, pp.~9 and~24]{Marker2014InfinitaryNotes}.
\end{hblemma}

The notes ask for the biconditional almost everywhere; with indexed
families it holds for every \(a\).

\begin{hbtheorem}{thm:p03-w03-p03-039-kueker-truth}{proved=annex:p03-w03-p03-039-kueker-truth}{Truth is approximated almost everywhere}
Let \(\tau\) be countable, \(\mathcal M\) a \(\tau\)-structure and
\(\varphi\) a sentence of \(\mathcal L_{\infty,\omega}(\tau)\), and let \(X\)
be large enough to approximate \(\mathcal M\) and \(\varphi\). Then
\(\mathcal M\models\varphi\) if and only if
\(\mathcal M^a\models\varphi^a\) a.e., that is, the set of \(a\) such that
\(\mathcal M^a\) is a substructure and \(\mathcal M^a\models\varphi^a\) is in
\(D(X)\) \cite[Thm.~2.43 and Ex.~2.44, pp.~24--25]{Marker2014InfinitaryNotes}.
\end{hbtheorem}

\begin{hbcorollary}{thm:p03-w03-p03-039-kueker-corollary}{proved=annex:p03-w03-p03-039-kueker-corollary}{Two special cases}
Under the hypotheses of \cref{thm:p03-w03-p03-039-kueker-truth}:
\begin{enumerate}
\item if \(\varphi\in\mathcal L_{\omega_1,\omega}\), then
      \(\mathcal M\models\varphi\) if and only if
      \(\mathcal M^a\models\varphi\) a.e.;
\item if \(\mathcal M\) is countable, then \(\mathcal M\models\varphi\) if
      and only if \(\mathcal M\models\varphi^a\) a.e.
\end{enumerate}
\cite[Cor.~2.45, p.~25]{Marker2014InfinitaryNotes}.
\end{hbcorollary}

\begin{hbtheorem}{thm:p03-w03-p03-039-kueker}{proved=annex:p03-w03-p03-039-kueker}{Kueker's characterization of \(\equiv_{\infty,\omega}\)}
Let \(\tau\) be countable and let \(X\) be large enough to approximate the
\(\tau\)-structures \(\mathcal M\) and \(\mathcal N\). Then
\begin{enumerate}
\item \(\mathcal M\equiv_{\infty,\omega}\mathcal N\) if and only if
      \(\mathcal M^a\cong\mathcal N^a\) a.e.;
\item \(\mathcal M\not\equiv_{\infty,\omega}\mathcal N\) if and only if
      \(\mathcal M^a\not\cong\mathcal N^a\) a.e.
\end{enumerate}
Here ``\(\mathcal M^a\cong\mathcal N^a\) a.e.'' means that the set of \(a\)
for which \(\mathcal M^a\) and \(\mathcal N^a\) are substructures and are
isomorphic is in \(D(X)\), and similarly for \(\not\cong\)
\cite[Thm.~2.46 and Ex.~2.47, pp.~25--26]{Marker2014InfinitaryNotes}.
\end{hbtheorem}

As the notes observe, \(D(X)\) is not an ultrafilter, so item~2 does not
follow from item~1 by complementation; each item needs its own forward
direction, and the backward directions follow because \(D(X)\) is a proper
filter. The notes obtain the forward direction of item~2 from
\cref{thm:p03-w03-p03-039-kueker-truth} applied to a sentence true in
\(\mathcal M\) and false in \(\mathcal N\), which needs \(X\) to approximate
that sentence as well; the annex argues with the relations \(\sim_\alpha\)
instead, so that \(X\) need only contain \(M\) and \(N\).

\begin{hbcorollary}{thm:p03-w03-p03-039-kueker-countable}{proved=annex:p03-w03-p03-039-kueker-countable}{Countable structures}
If \(\tau\) is countable, \(\mathcal M\) is countable and \(X\) is large
enough to approximate \(\mathcal M\) and \(\mathcal N\), then
\(\mathcal M\equiv_{\infty,\omega}\mathcal N\) if and only if
\(\mathcal M\cong\mathcal N^a\) a.e.
\cite[Cor.~2.48, p.~26]{Marker2014InfinitaryNotes}.
\end{hbcorollary}

\begin{hbcorollary}{thm:p03-w03-p03-039-kueker-free}{proved=annex:p03-w03-p03-039-kueker-free}{Free abelian groups through countable approximations}
Let \(F\) be the free abelian group of countably infinite rank and \(H\) an
abelian group. Then \(H\equiv_{\infty,\omega}F\) if and only if \(H\) is
\(\aleph_1\)-free and not free on a finite set of generators.
\end{hbcorollary}

The notes remark that \cref{thm:p03-w03-p03-039-kueker-countable} gives
another proof that a group is \(\equiv_{\infty,\omega}\)-equivalent to a
free abelian group exactly when its countable subgroups are free
\cite[p.~26]{Marker2014InfinitaryNotes}. Stated for the free group of
countably infinite rank, as above, the ``if'' direction needs the extra
hypothesis that \(H\) is not free on finitely many generators: the group
\(\mathbb Z\) is \(\aleph_1\)-free but satisfies
\(\exists x\,\forall y\bigvee_{n\in\mathbb Z}y=nx\), where \(nx\) is the
\(n\)-fold sum \(x+\cdots+x\) for \(n>0\), its inverse for \(n<0\) and \(0\)
for \(n=0\); \(F\) does not, as it is not cyclic.
The corollary recovers \cref{thm:p03-w03-p03-039-aleph-one-free} and the
second half of \cref{thm:p03-w03-p03-039-subgroups-free} by this second
route.

\subsection{Relation to the other constructions in this handbook}
\label{sec:p03-w03-p03-039-bridge}

\paragraph{Approximants.} \Cref{def:p03-morley-approximants} defines
\(\equiv_\alpha\) on pointed structures with the successor clause ``\(\equiv_\alpha\)
holds and the forth and back conditions hold''. The relation
\(\sim_\alpha\) of \cref{def:p03-w03-p03-039-approximants} omits the first
conjunct. The two relations coincide at every ordinal: they agree at
\(0\) and at limits, and at a successor the missing conjunct is implied by
the forth condition by \cref{thm:p03-w03-p03-039-approximant-basics}. The
symbol \(\equiv_\alpha\) of \cref{def:p03-w03-p03-039-approximants}, by
contrast, is agreement on sentences of quantifier rank at most \(\alpha\);
\cref{thm:p03-w03-p03-039-rank-characterization} shows that it coincides with
the unpointed \(\sim_\alpha\).

\paragraph{Scott formulas and sentences.} The formulas
\(\theta^\alpha_{\bar a}\) of \cref{def:p03-morley-scott-formulas} contain
the extra conjunct \(\theta^\alpha_{\bar a}\) at successor stages, while
\(\Phi^{\mathcal M}_{\bar a,\alpha+1}\) does not; both define the same
relation, by \cref{thm:p03-morley-scott-stabilization} and
\cref{thm:p03-w03-p03-039-scott-formula-lemma}. The Morley-route sentence
\(\sigma_M\) of \cref{thm:p03-morley-scott-sentence} is built at a countable
stabilization stage \(\delta_M\) for an at most countable \(M\) and is
proved to characterize \(M\) among at most countable structures. The
sentence \(\Phi^{\mathcal M}\) is built at the Scott rank
\(\operatorname{sr}(\mathcal M)\) for an infinite \(\mathcal M\) of any
cardinality, lies in \(\mathcal L_{\kappa^+,\omega}\), and characterizes the
\(\equiv_{\infty,\omega}\)-class of \(\mathcal M\) among all structures; the
countable characterization is then a consequence through
\cref{thm:p03-w03-p03-039-countable-back-forth}. The two proofs are
different routes: the Morley route builds an isomorphism directly by back
and forth on the formulas, while this route passes through Karp's partial
isomorphism systems.

\paragraph{Expressive power and complete sentences.} The axiomatizations of
\cref{thm:p03-w03-p03-039-axiomatizable} and the ordinal sentences used in
\cref{thm:p03-w03-p03-039-ordinal-equivalence} are the same sentences as
Nadel's (\cref{thm:p03-w03-p03-039-nadel-expressive} and
\cref{thm:p03-w03-p03-039-nadel-ordinals}); the notes and the chapter pose
them as exercises in different vocabularies (the notes' torsion groups are
abelian and written additively), and each is proved once. The results on
complete sentences in \cref{sec:p03-w03-p03-039-complete-sentences} all rest
on Scott's theorem in the form of \cref{thm:p03-w03-p03-039-scott}: a
sentence of \(\mathcal L_{\omega_1,\omega}\) over a countable vocabulary is
complete exactly when it implies the Scott sentence of one countable
structure, which is the content of
\cref{thm:p03-w03-p03-039-complete-categorical}. For a reader coming from
\cref{def:p02-lomega-one-omega}, nothing changes in passing to these results:
they concern sentences, and the two definitions of the logic have the same
sentences.

\paragraph{Atomic models and countable approximations.}
\Cref{thm:p03-w03-p03-039-complete-atomic} sharpens
\cref{thm:p03-w03-p03-039-chang} for complete sentences: Chang's reduction
needs a theory together with a countable family of types to omit, while for
a complete sentence the Scott formulas of the countable model, added to the
fragment, make every model's natural expansion atomic, so that the types are
omitted by atomicity alone. \Cref{thm:p03-w03-p03-039-kueker} and Karp's
\cref{thm:p03-w03-p03-039-karp} describe the same relation by different
means: Karp's partial isomorphism system is the object from which the club
of \cref{thm:p03-w03-p03-039-kueker} is built, and conversely the
isomorphisms \(\mathcal M^a\cong\mathcal N^a\) are back-and-forth systems
between countable pieces. \Cref{thm:p03-w03-p03-039-kueker-free} is the
second route, via countable approximations, to the group-theoretic
corollary proved first from Scott's theorem.

\paragraph{Basic calibration.} \Cref{thm:p02-lomega-one-omega-calibration}
records the failure of compactness and the existence of Scott sentences
for countable structures as a cited statement; those two claims are proved
here as \cref{thm:p03-w03-p03-039-compactness-failure} and
\cref{thm:p03-w03-p03-039-scott}.

\section{The Logic \texorpdfstring{\(\mathscr L_{\omega_1\omega}\)}{L omega1 omega} and Admissible Fragments after Nadel}
\label{sec:p03-w03-p03-039-nadel}

This section follows a second presentation, Nadel's survey chapter
``\(\mathscr L_{\omega_1\omega}\) and Admissible Fragments'' in the
Barwise--Feferman volume \emph{Model-Theoretic Logics}
\cite{Nadel1985Admissible}. The text read is the open-access copy of the
chapter in the Perspectives in Logic reissue of that volume, whose pages carry
the original pagination (pp.~271--316). Nadel's notation is kept:
\(\mathscr L_{\infty\omega}\), \(\mathscr L_{\kappa\omega}\) and
\(\mathscr L_{\omega_1\omega}\) for the logics, \(\mathfrak M,\mathfrak N\)
for structures with universes \(M,N\), \(L_B(\tau)\) for a fragment, and
\(S\), \(s\) for a consistency property and its members. The chapter is a
survey: most of its results are stated with a reference and no proof. Every
result recorded below is proved in the annex, and where the chapter's
printed statement needs a correction the correction is stated and justified
next to it. The purpose is to give the reader Nadel's route through the
subject: the logics and their expressive power, the back-and-forth
characterization of \(\equiv_{\infty\omega}\), the model existence theorem
for countable fragments, and the set-theoretic data (\(\Delta_0\)- and
\(\Sigma\)-formulas, the theory \(\mathrm{KP}\), admissible sets) from which
admissible fragments are cut. The comparison with the presentation of
\cref{sec:p03-w03-p03-039-marker} is made result by result.

\subsection{The infinitary logics and their expressive power}
\label{sec:p03-w03-p03-039-nadel-logics}

\begin{hbdefinition}{def:p03-w03-p03-039-nadel-logics}{The logics \(\mathscr L_{\infty\omega}\), \(\mathscr L_{\kappa\omega}\), \(\mathscr L_{\infty\infty}\) and \(\mathscr L_{\infty\lambda}\)}
The formulas of \(\mathscr L_{\infty\omega}(\tau)\) form the smallest class
that contains the atomic formulas of \(\tau\), is closed under the
connectives and the single-variable quantifiers of first-order logic, and
contains \(\bigwedge\Phi\) for every \emph{set} \(\Phi\) of its formulas.
\(\mathfrak M\models\bigwedge\Phi\) under an assignment when every member of
\(\Phi\) holds under it; \(\bigvee\Phi\) abbreviates
\(\neg\bigwedge\{\neg\varphi:\varphi\in\Phi\}\). From here on only formulas
with finitely many free variables are considered. For an infinite regular
cardinal \(\kappa\), \(\mathscr L_{\kappa\omega}\) is the sublogic in which
every conjunction is over a set of cardinality less than \(\kappa\); thus
\(\mathscr L_{\omega_1\omega}\) has countable conjunctions and disjunctions
and \(\mathscr L_{\omega\omega}\) is first-order logic. The logic
\(\mathscr L_{\infty\infty}\) adds the formula \(\exists X\varphi\) (and
\(\forall X\varphi\)) for every set \(X\) of variables;
\(\mathscr L_{\infty\lambda}\) is its sublogic in which quantified sets of
variables have fewer than \(\lambda\) members, and only formulas with fewer
than \(\lambda\) free variables are considered. The structures are ordinary
\(\tau\)-structures, and \(\equiv_{\infty\omega}\) is agreement on
\(\mathscr L_{\infty\omega}\)-sentences
\cite[\S\S~1.2 and~2, pp.~273 and~276]{Nadel1985Admissible}.
\end{hbdefinition}

For singular \(\kappa\) the chapter says that the definition of
\(\mathscr L_{\kappa\omega}\) differs, so that conjunctions of conjunctions
cannot simulate a conjunction of size \(\kappa\), and it does not give that
definition \cite[p.~273]{Nadel1985Admissible}; no logic
\(\mathscr L_{\kappa\omega}\) with \(\kappa\) singular is used below. The
logic \(\mathscr L_{\infty\omega}\) here has the same sentences as the logic
of \cref{def:p03-w03-p03-039-language}; the two definitions differ only in
taking \(\bigvee\) as primitive or defined and in the restriction to finitely
many free variables.

\begin{hblemma}{thm:p03-w03-p03-039-nadel-free-variables}{proved=annex:p03-w03-p03-039-nadel-free-variables}{Subformulas keep finitely many free variables}
If a formula \(\varphi\) of \(\mathscr L_{\infty\omega}\) has finitely many
free variables, then so does every subformula of \(\varphi\)
\cite[p.~273]{Nadel1985Admissible}.
\end{hblemma}

\begin{hbtheorem}{thm:p03-w03-p03-039-nadel-expressive}{proved=annex:p03-w03-p03-039-nadel-expressive}{Five properties expressible in \(\mathscr L_{\omega_1\omega}\)}
Each of the following classes is the class of models of one sentence of
\(\mathscr L_{\omega_1\omega}\) (in the vocabulary indicated):
\begin{enumerate}
\item the finite structures, in the empty vocabulary (equality only);
\item the torsion groups, in the vocabulary \(\{\cdot,{}^{-1},1\}\), among groups;
\item the finitely generated groups, in the same vocabulary, among groups;
\item the \(\omega\)-models: structures for a vocabulary with a unary
      relation symbol \(\mathsf N\) and constants \(\underline n\)
      \((n\in\omega)\) in which \(\mathsf N\) is interpreted as exactly the
      set of values of the constants and distinct constants have distinct
      values;
\item for a countable set \(\Phi(x_1,\dots,x_k)\) of formulas of
      \(\mathscr L_{\omega_1\omega}\) (in particular a type of
      \(\mathscr L_{\omega\omega}\)), the structures that omit \(\Phi\),
      that is, in which no tuple satisfies every member of \(\Phi\)
\end{enumerate}
\cite[\S\S~1.1 and~1.3, pp.~272--274]{Nadel1985Admissible}.
\end{hbtheorem}

\begin{hbtheorem}{thm:p03-w03-p03-039-nadel-well-orders}{proved=annex:p03-w03-p03-039-nadel-well-orders}{Well-orderings in \(\mathscr L_{\omega_1\omega_1}\)}
Let \(\sigma_{\mathrm{wo}}\) be the conjunction of the axioms of a strict
linear order for \(<\) with
\(\neg\exists\{v_n:n\in\omega\}\bigwedge_{n\in\omega}v_{n+1}<v_n\), a
sentence of \(\mathscr L_{\omega_1\omega_1}\). Assuming the axiom of choice,
a structure \((L,<)\) satisfies \(\sigma_{\mathrm{wo}}\) if and only if it is
a well-ordering \cite[p.~274]{Nadel1985Admissible}.
\end{hbtheorem}

The proof uses choice (dependent choices suffices) in one direction only: a
well-ordering has no infinite descending sequence without any choice, while
passing from ``not well-founded'' to an infinite descending sequence needs
it. This is the reason for the chapter's qualification that the
characterization holds at least under the axiom of choice
\cite[p.~274]{Nadel1985Admissible}.

\begin{hbtheorem}{thm:p03-w03-p03-039-nadel-ordinals}{proved=annex:p03-w03-p03-039-nadel-ordinals}{Each ordinal is characterized by a sentence of \(\mathscr L_{\infty\omega}\)}
In the vocabulary \(\{<\}\) define formulas \(\varphi_\beta(x)\) by recursion
on the ordinal \(\beta\):
\[
\varphi_\beta(x)\;=\;\forall y\Bigl(y<x\to\bigvee_{\gamma<\beta}\varphi_\gamma(y)\Bigr)\;\wedge\;\bigwedge_{\gamma<\beta}\exists y\bigl(y<x\wedge\varphi_\gamma(y)\bigr),
\]
and let \(\sigma_\alpha\) be the conjunction of the strict linear order axioms
with \(\forall x\bigvee_{\beta<\alpha}\varphi_\beta(x)\) and
\(\bigwedge_{\beta<\alpha}\exists x\,\varphi_\beta(x)\). For every linear order
\((L,<)\) and \(a\in L\), \((L,<)\models\varphi_\beta(a)\) if and only if the
initial segment \(\{y\in L:y<a\}\) is well-ordered of order type \(\beta\);
and for every ordinal \(\alpha\ge1\) and every \(\{<\}\)-structure
\(\mathfrak M\), \(\mathfrak M\models\sigma_\alpha\) if and only if
\(\mathfrak M\cong(\alpha,<)\). The chapter poses this as an exercise to be
done by induction on the ordinal \cite[p.~274]{Nadel1985Admissible}.
\end{hbtheorem}

Since structures are nonempty, the case \(\alpha=0\) does not arise. The
empty disjunction is false and the empty conjunction is true, so
\(\varphi_0(x)\) says that \(x\) has no predecessor.

\subsection{Elementary equivalence: back and forth}
\label{sec:p03-w03-p03-039-nadel-equivalence}

\begin{hbdefinition}{def:p03-w03-p03-039-nadel-back-forth}{Partial isomorphisms and \(\kappa\)-back-and-forth sets}
A function \(f\) from \(\mathfrak M\) to \(\mathfrak N\) (same vocabulary) is
a \emph{partial isomorphism} if \(f\) extends to an isomorphism of the
substructure of \(\mathfrak M\) generated by \(\operatorname{dom}f\) onto the
substructure of \(\mathfrak N\) generated by \(\operatorname{ra}f\). For a
cardinal \(\kappa\), a nonempty set \(F\) of partial isomorphisms from \(\mathfrak M\)
to \(\mathfrak N\) is a \emph{\(\kappa\)-back and forth set} if for every
\(f\in F\):
\begin{enumerate}
\item[(i)] \(\forall X\subseteq M\,[\,\lvert X\rvert<\kappa\to\exists g\in F\,(f\subseteq g\wedge X\subseteq\operatorname{dom}g)\,]\);
\item[(ii)] \(\forall Y\subseteq N\,[\,\lvert Y\rvert<\kappa\to\exists h\in F\,(f\subseteq h\wedge Y\subseteq\operatorname{ra}h)\,]\).
\end{enumerate}
\(\mathfrak M\cong_{p,\kappa}\mathfrak N\) (\(\kappa\)-partially isomorphic)
means that such a set exists, and \(\cong_p\) is \(\cong_{p,\omega}\). A logic
has the \emph{Karp property} if \(\mathfrak M\cong_p\mathfrak N\) implies
that \(\mathfrak M\) and \(\mathfrak N\) are elementarily equivalent in that
logic \cite[\S~2.1, pp.~276--277]{Nadel1985Admissible}.
\end{hbdefinition}

\begin{hbwarning}{def:p03-w03-p03-039-nadel-nonempty}{The back-and-forth set must be nonempty}
Without the nonemptiness condition, clauses (i) and (ii) would hold
vacuously for the empty set between any two structures. Under that reading,
\(\cong_{p,\kappa}\) would hold between a one-element and a two-element set
in the empty vocabulary, which disagree on the first-order sentence
\(\forall x\forall y\,x=y\). The definition above therefore includes
\(F\ne\varnothing\); this is the condition under which
\cref{thm:p03-w03-p03-039-nadel-karp} holds. The substructure generated by
the empty set is allowed to be empty, so the empty map is a partial
isomorphism exactly when the constant symbols generate isomorphic
substructures by the map sending each constant to the same constant.
\end{hbwarning}

\begin{hblemma}{thm:p03-w03-p03-039-nadel-bridge}{proved=annex:p03-w03-p03-039-nadel-bridge}{Nadel's and Marker's partial isomorphisms agree}
\begin{enumerate}
\item A map \(f\) with \(\operatorname{dom}f\subseteq M\) and values in
      \(N\) is a partial isomorphism in the sense of
      \cref{def:p03-w03-p03-039-nadel-back-forth} if and only if it is a
      partial \(\tau\)-embedding in the sense of
      \cref{def:p03-w03-p03-039-partial-iso}.
\item For \(2\le n<\omega\), a nonempty set \(F\) of partial isomorphisms is
      an \(n\)-back and forth set if and only if it is an \(\omega\)-back and
      forth set, if and only if it is a partial isomorphism system in the
      sense of \cref{def:p03-w03-p03-039-partial-iso}.
\end{enumerate}
In particular \(\mathfrak M\cong_p\mathfrak N\) (with a nonempty set) in the
sense of \cref{def:p03-w03-p03-039-nadel-back-forth} if and only if
\(\mathfrak M\cong_p\mathfrak N\) in the sense of
\cref{def:p03-w03-p03-039-partial-iso}. The equivalence of \(\kappa=\omega\)
with \(\kappa=n\) is stated in the chapter \cite[p.~276]{Nadel1985Admissible}.
\end{hblemma}

\begin{hbtheorem}{thm:p03-w03-p03-039-nadel-karp}{proved=annex:p03-w03-p03-039-nadel-karp}{Karp's theorem (Nadel's Theorem~2.1.1)}
\(\mathfrak M\equiv_{\infty\omega}\mathfrak N\) if and only if
\(\mathfrak M\cong_p\mathfrak N\)
\cite[Thm.~2.1.1, p.~277]{Nadel1985Admissible}.
\end{hbtheorem}

The chapter says the theorem is proved by a direct induction on formulas and
refers to its Chapter~IX for the details \cite[p.~277]{Nadel1985Admissible}.
The annex proof is a different route: it reduces the statement through
\cref{thm:p03-w03-p03-039-nadel-bridge} to
\cref{thm:p03-w03-p03-039-karp}, whose proof is the induction on formulas
and the Scott-formula construction of the first section.

\begin{hbcorollary}{thm:p03-w03-p03-039-nadel-countable}{proved=annex:p03-w03-p03-039-nadel-countable}{The weak form of Scott's theorem (Nadel's Corollary~2.1.2)}
If \(\lvert M\rvert=\lvert N\rvert=\aleph_0\) and
\(\mathfrak M\equiv_{\infty\omega}\mathfrak N\), then
\(\mathfrak M\cong\mathfrak N\)
\cite[Cor.~2.1.2, p.~277]{Nadel1985Admissible}.
\end{hbcorollary}

\begin{hbexample}{def:p03-w03-p03-039-nadel-examples}{Dense orders, algebraically closed fields and ordinals}
The chapter offers four test cases \cite[pp.~277--278]{Nadel1985Admissible}.
Fields are taken in the vocabulary \(\{+,-,\cdot,0,1\}\), and transcendence
rank means transcendence degree over the prime field.
\begin{enumerate}
\item Any two dense linear orderings without endpoints are partially
      isomorphic (\cref{thm:p03-w03-p03-039-nadel-dlo}).
\item Algebraically closed fields of distinct finite transcendence degrees
      are not partially isomorphic (\cref{thm:p03-w03-p03-039-nadel-acf-finite}).
\item Algebraically closed fields of infinite transcendence degree are
      \(\equiv_{\infty\omega}\) when they have the same characteristic
      (\cref{thm:p03-w03-p03-039-nadel-acf-infinite}).
\item Well-orderings of distinct order types are not \(\equiv_{\infty\omega}\)
      (\cref{thm:p03-w03-p03-039-nadel-ordinals-distinct}).
\end{enumerate}
In the chapter the last two items are drawn from the potential-isomorphism
characterization of its \S~2.2; the proofs below use only partial
isomorphisms and the sentences of
\cref{thm:p03-w03-p03-039-nadel-ordinals}.
\end{hbexample}

\begin{hbwarning}{def:p03-w03-p03-039-nadel-characteristic}{Characteristic in the field example}
The chapter states that any two algebraically closed fields of infinite
transcendence rank are \(\equiv_{\infty\omega}\)
\cite[p.~278]{Nadel1985Admissible}. Without the hypothesis of equal
characteristic this fails: an algebraic closure of
\(\mathbb Q(x_n:n\in\omega)\) satisfies \(\neg\,1+1=0\), and an algebraic
closure of \(\mathbb F_2(x_n:n\in\omega)\) satisfies \(1+1=0\). The
hypothesis is added in \cref{thm:p03-w03-p03-039-nadel-acf-infinite}.
\end{hbwarning}

\begin{hbconvention}{def:p03-w03-p03-039-nadel-field-facts}{Facts from field theory used below}
The field examples use four facts of algebra, applied in the vocabulary
\(\{+,-,\cdot,0,1\}\): (F1) if \(K\) is a subfield of an algebraically closed
field \(E\), the set \(\overline K^E\) of elements of \(E\) algebraic over
\(K\) is an algebraically closed field algebraic over \(K\), that is, an
algebraic closure of \(K\); (F2) an isomorphism of fields \(K\cong K'\)
extends to an isomorphism of any algebraic closure of \(K\) onto any
algebraic closure of \(K'\) (Steinitz); (F3) if \(a\) is transcendental over
\(K\) and \(b\) is transcendental over \(K'\), an isomorphism \(K\cong K'\)
extends to an isomorphism \(K(a)\cong K'(b)\) sending \(a\) to \(b\); (F4)
if \(L\) is algebraic over \(K\) and \(c\) is algebraic over \(L\), then \(c\)
is algebraic over \(K\). The
transcendence degree is the supremum of the cardinalities of sets
algebraically independent over the prime field. The distinction needed
below is between a finite bound on the sizes of independent tuples and
the existence of arbitrarily large finite independent tuples.
\end{hbconvention}

\begin{hblemma}{thm:p03-w03-p03-039-nadel-field-dependencies}{proved=annex:p03-w03-p03-039-nadel-field-dependencies}{Algebraic extensions and the finite-generator bound}
In ZFC the four field facts (F1)--(F4) of
\cref{def:p03-w03-p03-039-nadel-field-facts} hold. In addition:
\begin{enumerate}
\item if a field \(E\) is algebraic over \(K(b_1,\ldots,b_r)\),
      every subset of \(E\) algebraically independent over \(K\) has
      at most \(r\) elements;
\item consequently a field of infinite transcendence degree over \(K\)
      is not algebraic over the relative algebraic closure of any
      finitely generated extension \(K(b_1,\ldots,b_r)\) inside it.
\end{enumerate}
The isomorphism-extension assertion (F2) uses Zorn's lemma. The annex
also proves the exchange argument needed for the finite-generator bound;
that bound is not inferred from the number of symbols in the field
language.
\end{hblemma}

\begin{hbtheorem}{thm:p03-w03-p03-039-nadel-dlo}{proved=annex:p03-w03-p03-039-nadel-dlo}{Dense orders without endpoints}
If \(\mathfrak M=(M,<)\) and \(\mathfrak N=(N,<)\) are dense linear orderings
without endpoints, the set of order-preserving maps with finite domain
from \(M\) to \(N\) is a nonempty \(\omega\)-back and forth set; hence
\(\mathfrak M\cong_p\mathfrak N\)
\cite[p.~277]{Nadel1985Admissible}.
\end{hbtheorem}

\begin{hbtheorem}{thm:p03-w03-p03-039-nadel-acf-finite}{proved=annex:p03-w03-p03-039-nadel-acf-finite}{Distinct finite transcendence degrees}
If \(\mathfrak M\) and \(\mathfrak N\) are algebraically closed fields of
finite transcendence degrees \(m\ne n\), then
\(\mathfrak M\not\cong_p\mathfrak N\)
\cite[p.~277]{Nadel1985Admissible}.
\end{hbtheorem}

\begin{hbtheorem}{thm:p03-w03-p03-039-nadel-acf-infinite}{proved=annex:p03-w03-p03-039-nadel-acf-infinite}{Infinite transcendence degree}
If \(\mathfrak M\) and \(\mathfrak N\) are algebraically closed fields of the
same characteristic, both of infinite transcendence degree, then
\(\mathfrak M\cong_p\mathfrak N\) and so
\(\mathfrak M\equiv_{\infty\omega}\mathfrak N\)
\cite[p.~278]{Nadel1985Admissible}.
\end{hbtheorem}

\begin{hbcorollary}{thm:p03-w03-p03-039-nadel-ordinals-distinct}{proved=annex:p03-w03-p03-039-nadel-ordinals-distinct}{Distinct ordinals are distinguished}
If \(\alpha\ne\beta\) are nonzero ordinals, then
\((\alpha,<)\not\equiv_{\infty\omega}(\beta,<)\)
\cite[p.~278]{Nadel1985Admissible}.
\end{hbcorollary}

Together with \cref{thm:p03-w03-p03-039-nadel-countable}, the field example
shows why the countability hypothesis there cannot be removed: an
algebraically closed field of transcendence degree \(\aleph_0\) and one of
degree \(\aleph_1\), in the same characteristic, are
\(\equiv_{\infty\omega}\) by
\cref{thm:p03-w03-p03-039-nadel-acf-infinite} but have different
cardinalities.

\subsection{Fragments, consistency properties and model existence}
\label{sec:p03-w03-p03-039-nadel-model-existence}

\begin{hbdefinition}{def:p03-w03-p03-039-nadel-fragment}{Fragment (Nadel's Definition~3.1.1)}
A \emph{fragment} of \(\mathscr L_{\infty\omega}(\tau)\) is a set \(L_B(\tau)\)
of formulas and variables of \(\mathscr L_{\infty\omega}(\tau)\) such that
\begin{enumerate}
\item[(i)] \(\mathscr L_{\omega\omega}(\tau)\subseteq L_B(\tau)\);
\item[(ii)] if \(\varphi\in L_B(\tau)\), every subformula and variable of
      \(\varphi\) is in \(L_B(\tau)\);
\item[(iii)] if \(\varphi(v)\in L_B(\tau)\) and \(\sigma\) is a term of
      \(\tau\) all of whose variables lie in \(L_B(\tau)\), then
      \(\varphi(\sigma/v)\in L_B(\tau)\);
\item[(iv)] if \(\varphi,\psi\) and \(v\) are in \(L_B(\tau)\), so are
      \(\neg\varphi\), \(\exists v\varphi\), \(\forall v\varphi\),
      \(\varphi\,\&\,\psi\), \(\varphi\vee\psi\) and \({\sim}\varphi\),
\end{enumerate}
where \({\sim}\theta\) is \(\neg\theta\) for atomic \(\theta\),
\({\sim}(\neg\theta)\) is \(\theta\), \({\sim}(\bigwedge\Theta)\) is
\(\bigvee\{{\sim}\theta:\theta\in\Theta\}\), \({\sim}(\bigvee\Theta)\) is
\(\bigwedge\{{\sim}\theta:\theta\in\Theta\}\), \({\sim}(\exists v\varphi)\)
is \(\forall v\neg\varphi\) and \({\sim}(\forall v\varphi)\) is
\(\exists v\neg\varphi\) \cite[Def.~3.1.1, pp.~278--279]{Nadel1985Admissible}.
\end{hbdefinition}

This is a different definition from \cref{def:p03-w03-p03-039-fragment}:
Nadel's fragment contains the whole of first-order logic over every
variable, and his formal negation of a quantified formula carries
\(\neg\) rather than \({\sim}\) inside the quantifier. The two are kept
apart; the model existence theorem is proved below for Nadel's version
directly. In this subsection \(\bigvee\Theta\) is treated as a formula in its
own right, as the clauses for \({\sim}\) require.

\begin{hbdefinition}{def:p03-w03-p03-039-nadel-consistency}{Consistency property (Nadel's Definition~3.1.2)}
Let the fragment \(L_B\) contain a set of constant symbols
\(C=\{c_n:n\in\omega\}\). A \emph{consistency property} for \(L_B\) is a set
\(S\) of sets of sentences of \(L_B\) such that for each \(s\in S\):
\begin{enumerate}
\item[(C0)] \(\varnothing\in S\); if \(s\subseteq s'\in S\), then
      \(s\cup\{\varphi\}\in S\) for each \(\varphi\in s'\);
\item[(C1)] if \(\varphi\in s\), then \(\neg\varphi\notin s\);
\item[(C2)] if \(\neg\varphi\in s\), then \(s\cup\{{\sim}\varphi\}\in S\);
\item[(C3)] if \(\bigwedge\Phi\in s\), then \(s\cup\{\varphi\}\in S\) for all \(\varphi\in\Phi\);
\item[(C4)] if \(\forall v\varphi(v)\in s\), then \(s\cup\{\varphi(c)\}\in S\) for every \(c\in C\);
\item[(C5)] if \(\bigvee\Phi\in s\), then \(s\cup\{\varphi\}\in S\) for some \(\varphi\in\Phi\);
\item[(C6)] if \(\exists v\varphi(v)\in s\), then \(s\cup\{\varphi(c)\}\in S\) for some \(c\in C\);
\item[(C7)] for every basic term \(t\) and all \(c,d\in C\):
      (i) if \((c=d)\in s\), then \(s\cup\{d=c\}\in S\);
      (ii) if \(\varphi(t)\in s\) and \((c=t)\in s\), then \(s\cup\{\varphi(c)\}\in S\);
      (iii) for some \(e\in C\), \(s\cup\{e=t\}\in S\).
\end{enumerate}
Here a \emph{basic term} is a member of \(C\), a constant symbol of
\(\tau\), or a term \(F(c_{i_1},\dots,c_{i_n})\) with \(F\) an \(n\)-ary
function symbol and \(c_{i_1},\dots,c_{i_n}\in C\), and in (C7)(ii)
\(\varphi(v)\) is a formula of \(L_B\) with one free variable \(v\)
\cite[Def.~3.1.2, p.~279]{Nadel1985Admissible}. A structure is
\emph{canonical} if each of its elements interprets a constant of \(C\).
\end{hbdefinition}

Two points of the printed definition are changed. In (C4) the chapter
prints ``\(s\cup\{\varphi(c)\}\in C\)'', a misprint for \(\in S\). In (C7)
the chapter lets \(t\) range only over terms \(F(c_{i_1},\dots,c_{i_n})\)
\cite[p.~279]{Nadel1985Admissible}; the next warning shows that the theorem
fails under that reading, and \(t\) is therefore allowed to be a member of
\(C\) or a constant symbol as well.

\begin{hbwarning}{def:p03-w03-p03-039-nadel-c7}{Equations between the constants of \(C\)}
Let \(\tau\) consist of the constants \(c_n\) \((n\in\omega)\) only, so that
there is no term \(F(c_{i_1},\dots,c_{i_n})\). Let \(L_B\) be any countable
fragment of \(\mathscr L_{\infty\omega}(\tau)\) and
\(S=\{\varnothing,\{\neg\,c_0=c_0\}\}\). Clause (C0) holds: \(\varnothing\in S\), and if
\(s\subseteq s'\in S\) and \(\varphi\in s'\), then \(s'=\{\neg\,c_0=c_0\}\),
\(\varphi\) is \(\neg\,c_0=c_0\) and
\(s\cup\{\varphi\}=\{\neg\,c_0=c_0\}\in S\); (C1) holds because \(\neg\neg\,c_0=c_0\) is in no
member; (C2) holds because \({\sim}(c_0=c_0)\) is \(\neg\,c_0=c_0\);
(C3)--(C6) are vacuous; and (C7) as printed is vacuous. But
\(\{\neg\,c_0=c_0\}\) has no model. With (C7) read for every basic term,
(C7)(iii) with \(t=c_0\) requires a set \(\{\neg\,c_0=c_0,e=c_0\}\in S\), and
\(S\) is not a consistency property.
\end{hbwarning}

\begin{hbtheorem}{thm:p03-w03-p03-039-nadel-model-existence}{proved=annex:p03-w03-p03-039-nadel-model-existence}{Model existence theorem (Nadel's Theorem~3.1.3)}
\begin{enumerate}
\item[(i)] Let \(L_B\) be a countable fragment and \(S\) a consistency
      property for \(L_B\). For each \(s\in S\) there is a canonical model
      \(\mathfrak M\models\bigwedge s\).
\item[(ii)] (Extended version.) If in addition \(T\) is a set of sentences
      of \(L_B\) such that \(s\cup\{\varphi\}\in S\) for each \(s\in S\) and
      \(\varphi\in T\), then for each \(s\in S\) there is a canonical model
      of \(T\cup s\).
\end{enumerate}
\cite[Thm.~3.1.3, p.~279]{Nadel1985Admissible}.
\end{hbtheorem}

The chapter writes the conclusion of (ii) as a model of
``\(T\cup\{s\}\)''; since \(s\) is a set of sentences the meaning is
\(T\cup s\). The chapter attributes the theorem to Makkai and its
formulation to Barwise, and states it without proof
\cite[p.~279]{Nadel1985Admissible}; the annex gives the Henkin
construction that the chapter describes as building the canonical model in
\(\omega\) stages. Unlike \cref{thm:p03-w03-p03-039-extended-model-existence},
(ii) needs no nonemptiness hypothesis, because (C0) puts \(\varnothing\) into
\(S\). The countability of \(L_B\) is used to schedule the construction in
\(\omega\) steps; the chapter remarks that the theorem fails without it
\cite[p.~280]{Nadel1985Admissible}.

\subsection{Admissible sets}
\label{sec:p03-w03-p03-039-nadel-admissible}

Syntax and satisfaction of \(\mathscr L_{\infty\omega}\) are taken to be
formalized in set theory, so that formulas and structures are sets. For a
transitive set \(B\), \(L_B=\mathscr L_{\infty\omega}\cap B\) is the set of
formulas that are members of \(B\) \cite[p.~288]{Nadel1985Admissible}. The
chapter's Part~II is about the sets \(B\) for which \(L_B\) is well behaved.

\begin{hbdefinition}{def:p03-w03-p03-039-nadel-levy}{\(\Delta_0\)-, \(\Sigma_1\)- and \(\Sigma\)-formulas; end extensions}
In a first-order language with a binary relation symbol \(\in\) (and possibly
other symbols), the \emph{\(\Delta_0\)-formulas} form the smallest class that
contains the atomic formulas and is closed under \(\neg\), \(\&\), \(\vee\)
and bounded quantification \(\exists u\in v\,\varphi\) (that is,
\(\exists u[u\in v\,\&\,\varphi]\)) and \(\forall u\in v\,\varphi\). The
\emph{\(\Sigma_1\)-formulas} are those of the form \(\exists v\varphi\) with
\(\varphi\) a \(\Delta_0\)-formula. The \emph{\(\Sigma\)-formulas} form the
closure of the \(\Delta_0\)-formulas under \(\&\), \(\vee\), bounded
quantifiers and existential quantifiers. A relation on a structure is
\(\Sigma\)-definable if a \(\Sigma\)-formula defines it without parameters,
\(\Sigma\) if a \(\Sigma\)-formula defines it with parameters, \(\Pi\) if
its complement is \(\Sigma\), and \(\Delta\) if it is both \(\Sigma\) and
\(\Pi\). A structure \(\langle B,F,\dots\rangle\) is an \emph{end extension}
of \(\langle A,E,\dots\rangle\), with \(E,F\) binary, if
\(\langle A,E,\dots\rangle\) is a substructure of \(\langle B,F,\dots\rangle\)
and \((c,a)\in F\) with \(a\in A\) implies \(c\in A\). A formula is
\emph{persistent} if, whenever it holds in a structure at a tuple of its
elements, it holds at that tuple in every end extension; it is
\emph{absolute} if it holds in the structure at the tuple exactly when it
holds there in every end extension. Relative to a theory \(T\) the same
notions are defined with all structures required to be models of \(T\)
\cite[pp.~288--289]{Nadel1985Admissible}.
\end{hbdefinition}

\begin{hblemma}{thm:p03-w03-p03-039-nadel-persistence}{proved=annex:p03-w03-p03-039-nadel-persistence}{\(\Delta_0\)-formulas are absolute and \(\Sigma\)-formulas persistent}
Let \(\langle B,F,\dots\rangle\) be an end extension of
\(\langle A,E,\dots\rangle\). Every \(\Delta_0\)-formula is absolute and
every \(\Sigma\)-formula is persistent between them. Relative to any theory
\(T\), a formula \(\varphi\) is absolute if and only if \(\varphi\) and
\(\neg\varphi\) are both persistent
\cite[p.~289]{Nadel1985Admissible}.
\end{hblemma}

The chapter also states the converse due to Feferman and Kreisel, that a
formula persistent relative to \(T\) is \(T\)-provably equivalent to a
\(\Sigma\)-formula, with a reference and no proof
\cite[p.~289]{Nadel1985Admissible}; it is not used below.

\begin{hbdefinition}{def:p03-w03-p03-039-nadel-kp}{The theory \(\mathrm{KP}\); admissible sets and structures}
\(\mathrm{KP}\) consists of the axioms of extensionality, pairing and union,
the foundation scheme for arbitrary formulas, and the two schemes
\begin{itemize}
\item \(\Delta_0\)-Separation:
      \(\exists v\,\forall x\,(x\in v\leftrightarrow x\in u\,\&\,\varphi(x))\),
      for each \(\Delta_0\)-formula \(\varphi\) in which \(v\) does not
      occur free;
\item \(\Delta_0\)-Collection:
      \(\forall x\in u\,\exists y\,\varphi(x,y)\to\exists v\,\forall x\in u\,\exists y\in v\,\varphi(x,y)\),
      for each \(\Delta_0\)-formula \(\varphi\) in which \(v\) does not
      occur free,
\end{itemize}
the remaining free variables of \(\varphi\) acting as parameters. A
structure \(\mathfrak A=\langle A,\in,\dots\rangle\) is \emph{admissible} if
\(\langle A,\in\rangle\) is transitive and
\(\langle A,\in,\dots\rangle\models\mathrm{KP}\), the schemes being taken for
the \(\Delta_0\)-formulas of the expanded vocabulary. A transitive set \(A\)
is \emph{admissible} if \(\langle A,\in\rangle\) is an admissible structure;
\(L_A\) is then the \emph{admissible fragment}. A transitive set satisfying
every axiom of \(\mathrm{KP}\) except perhaps \(\Delta_0\)-Collection is
\emph{rudimentary}. For a transitive \(B\), \(o(B)\) is the least ordinal not
in \(B\) \cite[pp.~289--290]{Nadel1985Admissible}.
\end{hbdefinition}

\begin{hbdefinition}{def:p03-w03-p03-039-nadel-hkappa}{Hereditarily small sets}
\(\mathrm{TC}(x)\) is the smallest transitive set having \(x\) as a subset.
For an infinite cardinal \(\kappa\),
\(H(\kappa)=\{x:\lvert\mathrm{TC}(x)\rvert<\kappa\}\);
\(\mathrm{HF}=H(\aleph_0)\) and \(\mathrm{HC}=H(\aleph_1)\)
\cite[p.~290]{Nadel1985Admissible}.
\end{hbdefinition}

\begin{hbtheorem}{thm:p03-w03-p03-039-nadel-hkappa}{proved=annex:p03-w03-p03-039-nadel-hkappa}{\(H(\kappa)\) is admissible for regular \(\kappa\)}
Assume the axiom of choice. For every infinite regular cardinal \(\kappa\),
\(H(\kappa)\) is an admissible set; in particular \(\mathrm{HF}\) and
\(\mathrm{HC}\) are admissible \cite[p.~290]{Nadel1985Admissible}.
\end{hbtheorem}

\(\mathrm{HF}\) is an admissible set with \(o(\mathrm{HF})=\omega\), which is
possible because \(\mathrm{KP}\) as defined above has no axiom of infinity.
The equation \(o(\mathrm{HF})=\omega\) is part of the annex proof.
